\documentclass[10pt,a4paper]{ltjsarticle}

\usepackage{amsmath,amssymb,amsthm}
\usepackage{enumitem}
\usepackage{booktabs}
\usepackage{array}
\usepackage[hidelinks]{hyperref}
\emergencystretch=2em

\newtheorem{theorem}{定理}[section]
\newtheorem{lemma}[theorem]{補題}
\newtheorem{proposition}[theorem]{命題}
\newtheorem{corollary}[theorem]{系}
\theoremstyle{definition}
\newtheorem{definition}[theorem]{定義}
\theoremstyle{remark}
\newtheorem{remark}[theorem]{注意}
\renewcommand{\proofname}{証明}

\newcommand{\cmp}{\operatorname{cmp}}
\newcommand{\obs}{\operatorname{obs}}
\newcommand{\qcmp}{\operatorname{qcmp}}
\newcommand{\D}{\mathcal D}
\newcommand{\SCL}{\operatorname{SCL}}
\newcommand{\Obj}{\mathsf O}
\newcommand{\Om}{\Omega}
\newcommand{\blank}{\square}
\newcommand{\relto}{\mathrel{\rightsquigarrow}}

\title{統語概念束における有限モノイド圧縮：\\アリティ階層と擬多様体の三分法}
\author{Takayuki Kuriyama\\
Independent Researcher, Tokyo, Japan\\
\texttt{growup.kuriyama@gmail.com}}
\date{}

\begin{document}
\maketitle

\begin{abstract}
Clark の統語概念束（syntactic concept lattice; SCL）は、言語の両側分布構造を記録するものであり、Wurm はこれを任意の有限アリティのタプルへ拡張した。
本稿では、主層（principal layer）上でアリティ $f$ までのガード付きタプル代入を保存する有限モノイド観測の像の最小サイズ
\(\cmp_f(L)\)
を研究する。

正則言語について、\(\cmp_f(L)\) を、点付き統語モノイドから出る $f$-分離的な関係的射（relational morphism）の終域の最小濃度として厳密に特徴づける。
擬多様体 \(\mathbf V\) に属するすべての言語について、これらの圧縮数が一様に安定する最小アリティを \(\operatorname{ch}(\mathbf V)\) と書く。
本稿の主結果は、一様高さに関する次の三分法である：
\[
\operatorname{ch}(\mathbf V)\in\{1,2,\infty\},
\qquad
\operatorname{ch}(\mathbf V)=\infty
\Longleftrightarrow
\operatorname{Synt}(\{ab\})\in\mathbf V.
\]
したがって、有限の一様圧縮高さ $3,4,\ldots$ は存在しない。

無限の場合の結果は鋭い。
\(\langle\operatorname{Synt}(\{ab\})\rangle\) の内部では、任意の境界 $d\to d+1$ において圧縮ギャップを非有界にでき、アリティ階層の任意の有限な真の初期部分を実現できる。
一方、有限側では、可換モノイドとバンドはアリティ1で安定し、完全正則な統語モノイドはすべてアリティ2までに安定する。有限群の核言語は、この二項の上界が鋭いことを示す。

単項アリティでは、任意の非空有限単純グラフが明示的な長さ3の言語によって実現される。
これにより彩色数による厳密な公式が得られ、明示的に列挙された長さ3の言語に対して \(\cmp_1(L)\le3\) を判定する問題が NP 完全であることが従う。
圧縮高さ1と2の構造的境界は未解決である。
\end{abstract}

\noindent\textbf{キーワード：} 統語概念束；有限モノイド；関係的射；擬多様体；タプル・アリティ；グラフ彩色。

\noindent\textbf{2020 Mathematics Subject Classification:}
20M35 (primary); 20M07, 68Q70, 68Q45, 68Q17.

\section{序論}

\subsection{有限合成的圧縮}

正則言語では、厳密な両側文脈挙動は、その言語を認識する標準的な有限モノイドである統語モノイド $T_L$ によって捉えられる。
このモノイドの大きさは標準的な代数的記述量であり、統語モノイドの大きさとオートマトンの大きさとの関係は、正則言語の記述複雑性の文脈で体系的に研究されてきた~\cite{HolzerKonig2004}。
本稿は、これより選択的な問いから出発する。
厳密な認識そのものではなく、ひとつの固定された合成的タスク---すなわち、受理文脈への語およびタプルのガード付き代入---だけを要求するなら、$T_L$ の情報を実際にはどこまで保持する必要があるのか。
したがって、以下で扱う不変量は、認識不変量そのものに取って代わるものではなく、統語的複雑性をタスク相対的に緩和した量とみなすのが適切である。
本稿を通じて「圧縮」とは、この有限観測の終域を小さくすることを意味し、直線的プログラムなどによる語や半群式そのものの圧縮を意味しない。

このような緩和が単なる商問題にはならないと期待すべき古典的な先例がある。
不完全指定有限状態機械の最小化では、状態の適合性は推移的であるとは限らず、そのため標準的な縮約は通常の同値類最小化ではなく compatible cover を用いて行われる。
Paull と Unger は最大 compatible／最小 closed cover の古典的枠組みを発展させた~\cite{PaullUnger1959}。
また、不完全指定機械の状態縮約は NP 完全である~\cite{Pfleeger1973}。
ここでの類似は構造的なものであり、両問題を同一視するものではない。
本稿の有限対象は遷移 cover ではなくモノイドであり、さらに積に関する合成可能性が、古典的な機械最小化問題にはない制約を課す。
しかし同じ障害パターンが再び現れる。
部分的な仕様からは、反射的かつ対称的だが推移的とは限らない適合関係が自然に生じるため、商による同一視は硬直的すぎる場合がある。
最小分離 DFA も、合成的検証において補完的なアルゴリズム上の先例を与える。そこでは、separator の大きさが再利用可能な有限仮定の大きさを支配する~\cite{ChenFarzanClarkeTsayWang2009,VazquezPargaGarciaLopez2016}。

Clark--Wurm の統語概念構成は、ここで不足していた要素を与える。すなわち、部分仕様を\emph{言語そのものの内部から}生成する。
Clark の統語概念束 $\SCL_1(L)$ は、個々の語に対する統語代数を、語とその受理両側文脈との Galois incidence へ持ち上げる~\cite{ClarkSCL}。
Wurm はこの構成を任意の有限アリティのタプルへ拡張した~\cite{WurmSCLTuples}。
語およびタプルは、それぞれの主概念（principal concept）を通じてこれらの束に入る。
異なる二つの主概念が真の代入衝突を生じるのは、ちょうどその受理 intent が重なるときである。
すなわち、両者はある共通の受理文脈ではともに出現できる一方、別の文脈では区別される。
定義~\ref{scl:def:safety} では、この幾何学的述語を一度だけ定義する。
そして Safety Interface Theorem（定理~\ref{alg:thm:safety-interface}）により、その代表元レベルの principal-SCL 形式、ガード付き語／文脈形式、合同形式、および標準的な関係的射形式がすべて同値であることを示す。
続く selector lemma により、この標準的関係インターフェースは、任意の関係的射上の最適化へ拡張される。

これにより、厳密な統語情報から Clark--Wurm の主衝突へ、さらにそこから最小の有限合成的 separator へ至る概念的な流れが得られる。
中間段階が、仕様を外部から与えるのではなく内在的なものにしている。
最後の段階はタスク相対的である。
本稿では概念束全体を保存しようとはせず、実際の語およびタプルを代入するために必要な区別だけを保存する。
この意味で「標準的（canonical）」という語は Clark--Wurm 幾何から生成される衝突仕様を指すのであり、すべての閉概念に対して唯一の圧縮概念が存在すると主張するものではない。

正則言語では、こうして得られる制約を厳密に有限代数へ還元できる。
アリティ1では、singleton から生成される概念の等しさは統語合同と一致する~\cite[Proposition~3.19]{ClarkSCL} ので、単項主層は標準的に $T_L$ である。
より高い SCL 階層は追加の分離制約を生成し、その一方で単項主モノイドが再利用可能な合成的パレットを与える。
アリティ $f$ まで同時に必要となるモノイド値の個数を最小化し、その最適値を\emph{有限モノイド圧縮数} $\cmp_f(L)$ と呼ぶ。
最適な観測は統語類上で一定である必要がないため、適切な有限代数意味論は一般には通常の商ではなく、点付き統語モノイドからの関係的射となる。
これは上記の cover 型現象の乗法的対応物であり、不完全指定有限状態機械（ISFSM）の最小化を本問題へ還元しているわけではない。

厳密な認識写像は常にひとつの witness になるので、
\[
\cmp_1(L)\le \cmp_2(L)\le\cdots\le |T_L|.
\]
この鎖は、要求する代入アリティを上げたとき、厳密な統語情報のうちどれだけが残らなければならないかを定量化する。
すでにアリティ1で非自明である。
$L=\{a,a^3\}$ では、文脈衝突グラフは2彩色可能であるにもかかわらず、連接と両立する2元モノイド彩色は存在せず、$\cmp_1(L)=3$ となる。
さらに鋭く、Section~\ref{sec:conflict} では単項族 $K_n=\{a^n,a^{n+2}\}$, $n\ge2$ を与え、
\[
\begin{aligned}
 \chi(\Gamma_{K_n})=2<3=\cmp_f(K_n)&<n+2=\qcmp_f(K_n)\\
 &<|T_{K_n}|=n+4 \qquad(f\ge1)
\end{aligned}
\]
を得る。
したがって、任意の有限圧縮から統語モノイドの商へ制限すると、必要状態数に非有界な損失が生じうる。
多値の関係的インターフェースは単なる技術的言い換えではない。
反対側の構造的極端として、群核言語 $K_\eta=\eta^{-1}(1_G)$ について後に
$\cmp_1=1$ および $f\ge2$ で $\cmp_f=[G:Z(G)]$ を得る。
さらに、任意の非空有限単純グラフは、明示的な長さ3の言語によって単項レベルで正確に実現される。
これらの例から、$\cmp_f$ は統語モノイドの大きさでも、SCL の大きさでも、通常のグラフ彩色数でもないことが分かる。
これは、SCL が生成する代入仕様を保存するために必要な最小の\emph{合成的}有限資源を測る量である。

\subsection{アリティ階層と主定理}

中心となる問いは、より高いタプル・アリティを要求したとき、必要な最小合成資源がどのように変化するかである。
\[
M_{ab}:=\operatorname{Synt}(\{ab\}),
\qquad
\mathbf V_{ab}:=\langle M_{ab}\rangle,
\]
と書く。ただし $M_{ab}=\{1,a,b,ab,0\}$ である。
個々の言語のレベルでは、階層は任意に多くの情報を保ちうる。
この一つの5元モノイドによって生成される非周期的擬多様体の内部だけで、任意の境界 $d\to d+1$ において圧縮の乗法的増加を非有界にできる。
さらに強く、任意の有限 $m$ に対して一つの言語で
\[
\cmp_1(L)<\cmp_2(L)<\cdots<\cmp_m(L)
\]
を実現でき、しかも連続する比を、あらかじめ与えられた任意の有限閾値より同時に大きくできる。
三角形的貼り合わせ原理（triangular gluing principle）が、この有限初期部分の普遍性の背後にある機構を切り出す。

ところが、擬多様体上で一様化すると、驚くべきことに逆の現象が起きる。
任意に長くなりうる局所階層が、三つの値しか取らない大域的スペクトルへ崩壊する。
$\operatorname{ch}(\mathbf V)$ を、統語モノイドが $\mathbf V$ に属するすべての正則言語について $\cmp_f$ が一様に安定し始める最小アリティとする。
本稿の主要な構造定理は
\[
 \boxed{\operatorname{ch}(\mathbf V)\in\{1,2,\infty\}}
\]
であり、さらに厳密に
\[
\operatorname{ch}(\mathbf V)=\infty
\Longleftrightarrow
M_{ab}\in\mathbf V
\Longleftrightarrow
\mathbf V_{ab}\subseteq\mathbf V
\]
と特徴づけられる。
したがって、有限の一様高さ $3,4,\ldots$ は起こらない。

証明では、同じ5元の障害物によって無限側と有限側を分離する。
$M_{ab}\in\mathbf V$ ならば、$\mathbf V_{ab}$ 内での全境界増幅が高さを無限にする。
逆に $M_{ab}\notin\mathbf V$ ならば、divisor による議論から
\[
\mathbf V\subseteq\mathbf{Com}
\qquad\text{または}\qquad
\mathbf V\subseteq\mathbf{CR}
\]
が従う。
可換な統語構造ではすでにアリティ1で階層が崩壊する。
有限側の主要な技術的入力はさらに強い。
完全正則な統語モノイドをもつ任意の正則言語について、二項分離的な関係的射はすべての有限アリティで分離的であり、したがって
\[
\cmp_2=\cmp_3=\cdots
\]
となる。
有限バンドはこれとは別の非可換な高さ1の領域を与え、有限群の核言語は二項の上界が鋭いことを示す。
すなわち $K_\eta=\eta^{-1}(1_G)$ に対して
\[
\cmp_1(K_\eta)=1,
\qquad
\cmp_f(K_\eta)=[G:Z(G)]\quad(f\ge2).
\]
高さ1と2の一般的な境界は未解決である。

単項最適化には、直接的な計算量論的側面もある。
任意の非空有限単純グラフ $G$ は、明示的な長さ3の言語 $L_G$ により
\[
\cmp_1(L_G)=\chi(G)
\]
を満たすように実現され、しかも余分な非孤立衝突頂点は生じない。
したがって、明示的に列挙された長さ3の言語について $\cmp_1(L)\le3$ を判定する問題は NP 完全である。
明示的に与えられた有限点付き統語モノイドに対する対応問題も NP 完全である。
このように同じ不変量が、任意の局所衝突幾何、強いアリティ増幅、そして擬多様体レベルでは剛直な安定化スペクトルという三つの顔をもつ。

\subsection{関連研究と本稿の寄与の範囲}

Clark は統語概念束を分布代数として導入し、統語合同および universal automaton との関係を示し、その最小／終対象的な認識役割を確立した~\cite{ClarkSCL}。
Wurm はこの構成を有限および無限のタプル・アリティへ拡張し、有限アリティにおける有限性基準を証明するとともに、無限タプル SCL の有限性を特徴づけた~\cite{WurmSCLTuples}。
関連するオートマトン意味論は~\cite{WurmJLLI2017} に現れ、Kuznetsov は無限項 action logic に対する最近の SCL 意味論を展開している~\cite{Kuznetsov2026}。
Clark--Yoshinaka は、多重文脈自由文法へ向けた別の SCL 型代数拡張を与えている~\cite{ClarkYoshinaka2014}。
著者の知る限り、これらの研究はいずれも、複数のタプル・アリティにわたって overlap-sensitive な主代入衝突を保存する単一の有限モノイド資源を最小化してはいない。
なお、タプル SCL そのものを本稿の新規性として主張するものではない。

本稿の最適化問題は、いくつかの古典的最小化理論の中間にも位置する。
統語モノイドの大きさは確立された代数的記述量である~\cite{HolzerKonig2004}。
本稿の鎖
$\cmp_1(L)\le\cmp_2(L)\le\cdots\le |T_L|$
は、ガード付き代入だけを保存すればよいとしたとき、この厳密な資源をどこまで圧縮できるかを問う。
より抽象的には、非推移的な適合関係のもとでの最小化は、不完全指定有限状態機械の古典的問題である。
Paull--Unger は maximal compatible と closed cover を導入し~\cite{PaullUnger1959}、対応する状態縮約問題は NP 完全である~\cite{Pfleeger1973}。
本稿ではこれを構造的な先例としてのみ用いる。
彼らの compatible cover は本稿の関係的射ではない。
後者はモノイド積を保存しなければならず、また本稿の適合制約は、不完全な機械仕様として外部から与えられるのではなく、言語そのものから生成される。

最小 consistent automaton および separating automaton は、有限 separator の大きさを最適化する第二の、アルゴリズム的先例を与える~\cite{PittWarmuth1993,ChenFarzanClarkeTsayWang2009,VazquezPargaGarciaLopez2016,LaumenSnelVaandrager2026}。
合成的検証では、小さな separating DFA が再利用可能な有限仮定として働く。
正則推論には、直接のグラフ彩色定式化も存在する~\cite{CostaFlorencioVerwer2014}。
本稿の最適化は二点で異なる。
第一に、分離仕様が主 intent の重なりから内部的に生成される。
第二に、separator はタプル・アリティをまたいで再利用できるモノイド積を備えなければならない。
例 $L=\{a,a^3\}$ は第二の差を明瞭にする。
その単項衝突グラフは2彩色可能だが、$\cmp_1(L)=3$ である。

関係的射は有限半群論における古典的な道具である。
本稿では、圧縮写像が統語類上で一定である必要がないため、これが自然かつ標準的に現れる。
したがって適切な有限意味論は、一般には商写像ではなく、統語モノイドから圧縮終域への関係である。
この fiber による見方は、pointlike 理論における関係的射の古典的利用と隣接している。
そこでは、fiber への同時所属が、指定された有限 target のクラスでは分離できない source 元を記録する~\cite{Almeida1994,RhodesSteinberg2009}。
しかし本稿が解いているのは pointlike-set 問題ではない。
本稿の関係は、一つの言語に付随する点付き unsafe tuple によって制約され、目的は\emph{単一の分離終域の最小濃度}を求めることである。
さらに後の命題~\ref{conf:prop:strict-ladder-family} は、この関係的自由度が量的にも本質的であることを示す。
統語商だけに制限すると、単項の有限言語でさえ非有界な加法的コストが生じる。

擬多様体レベルの不変量 $\operatorname{ch}(\mathbf V)$ には、これとは別の方法論的先例がある。
代数的複雑性理論では、有限モノイドのクラス上で一様にプログラムを研究する~\cite{BarringtonTherien1988}。
より最近の研究では、特定のモノイドクラス内部の tame 性や資源階層が同定されている~\cite{GrosshansMcKenzieSegoufin2017}。
ここで技術的な還元を主張するものではない。
本稿の資源は終域濃度であり、パラメータはタプル・アリティである。
共通するのは、局所的な資源階層を、代数的に定義された source class 上で一様化するという点である。

ガード付き代入条件は分布的文法推論に由来する。
substitutable CFG については Clark--Eyraud を、多次元 substitutability については Yoshinaka を参照されたい~\cite{ClarkEyraud2007,Yoshinaka2009,Yoshinaka2011}。
著者の companion MCFG manuscript~\cite{KuriyamaMCFG2026} では、同じ数値的最小値を $\obs_f$ と書き、学習理論的な情報パラメータとして用いている。
したがって、この数値的不変量そのものを本稿の新規性として主張するものではない。
本稿の新しい寄与は、数値的最小値そのものではなく、その内在的な有限半群構造である。
具体的には、正則言語に対する厳密な関係的射による特徴づけ、そこから得られるアリティ階層と増幅現象、完全正則な場合の飽和定理、そして擬多様体レベルの圧縮高さ三分法である。
companion manuscript の学習、再構成、文法 rank に関する結果は、以下のいかなる証明にも用いない。

最後に、擬多様体および Rees 行列に関する背景は標準的である~\cite{Pin1986,Almeida1994,RhodesSteinberg2009,Howie1995,Petrich1969,Petrich1974}。
また $\operatorname{Synt}(\{ab\})$ は、古典的な最小非可換非周期生成子の一つとして現れる~\cite{MargolisPin1984}。
compatible tolerance も古典的である~\cite{ZelinkaTolerance1975,Pondelicek1978,Loganathan1986,KumaresanTolerance1984}。
現代的な用語では、Barber--Ru\v{s}kuc は反射的 compatible relation を diagonal subsemigroup として研究し、Rees 行列の合同構造を適応してそれらを記述している~\cite{BarberRuskucDirectSquare,BarberRuskucDiagonal2026}。
しかし、これらの一般結果だけから Section~\ref{sec:cr-saturation} の飽和定理が得られるわけではない。
本稿の fiber-overlap relation は特殊な対称 diagonal subsemigroup であるが、さらに $P$ に相対的な\emph{点付き}二項分離条件を課すことで、row/column の剛性と一意な乗法的 central defect が強制される。
新しい結論は全アリティ安定化定理であって、diagonal subsemigroup の分類ではない。

\subsection{論文の構成}

Sections~\ref{sec:scl-geometry}--\ref{sec:algebraic-conflict} では、SCL、語／文脈、合同、関係的射という各インターフェースを構成し、最後に可換な場合の高さ1の上界を得る。
Section~\ref{sec:vab-arity} では全境界増幅と有限初期部分の普遍性を証明し、圧縮高さを導入する。
Section~\ref{sec:cr-saturation} では完全正則モノイドに対するアリティ2飽和定理を証明する。
Section~\ref{sec:global-arity-trichotomy} では、これらと構造的 divisor argument を組み合わせて大域三分法を得るとともに、群・バンド・標準例における鋭い場合を記録する。
Section~\ref{sec:graph-universality} ではグラフ実現と NP 完全性を与え、Section~\ref{sec:discussion} では残された問題を述べて論文を閉じる。

\section{統語概念幾何と有限モノイド圧縮}
\label{sec:scl-geometry}

\paragraph{記号一覧。}
頻繁に用いる記号を参照用にまとめておく。
\begin{center}
\small
\begin{tabular}{@{}l p{0.72\linewidth}@{}}
\toprule
記号 & 意味\\
\midrule
\(\D_L^{(d)}(\vec x)\) & \(\vec x\) を受理する $d$-タプル文脈\\
\(\Obj_d(L),\Om_d(L)\) & 主 $d$ 項 SCL 対象概念とその重なりグラフ\\
\(\Gamma_L,\Gamma_L^{\mathrm{syn}}\) & 語レベルおよび統語的な単項衝突グラフ\\
\(\mathfrak U_d(L),\mathcal U_d(T,P)\) & 意味論的および有限統語的 unsafe-pair 関係\\
\(\Phi_d(\mathbf t)\) & $d$-タプルの有限統語 profile\\
\(\operatorname{Safe}_f(h;L)\) & アリティ $f$ までの圧縮写像の安全性\\
\(\cmp_f(L),\qcmp_f(L)\) & 無制限および商制限付き有限モノイド圧縮数\\
\(\operatorname{ch}(\mathbf V)\) & 擬多様体の一様 SCL 圧縮高さ\\
\bottomrule
\end{tabular}
\end{center}

\subsection{Clark--Wurm の有限タプル塔}

$L\subseteq\Sigma^*$ および $d\ge1$ を固定する。
Wurm~\cite{WurmSCLTuples} に従い、アリティ $d$ の対象集合を $(\Sigma^*)^d$、文脈集合を $(\Sigma^*)^{d+1}$ とし、incidence relation を
\[
(w_1,\ldots,w_d)\ I_d\ (x_0,\ldots,x_d)
\quad\Longleftrightarrow\quad
x_0w_1x_1\cdots w_dx_d\in L
\]
で定める。
これと同値な固定順序の hole 記法
\[
E=x_0\blank_1x_1\blank_2\cdots x_{d-1}\blank_dx_d
\]
を用いる。このとき
$E[\vec w]=x_0w_1x_1\cdots w_dx_d$ である。
したがって本稿で「タプル文脈」と言うときは常に、この順序を保存する Clark--Wurm 文脈を意味する。
タプルの各成分には空語を許す。
置換によって順序を変えるような拡張文脈概念は、本稿では一切用いない。
$A\subseteq(\Sigma^*)^d$ および $C\subseteq(\Sigma^*)^{d+1}$ に対して
\[
A^\uparrow=\{\mathbf x:\mathbf w I_d\mathbf x\text{ がすべての }\mathbf w\in A\text{ について成り立つ}\},
\qquad
C^\downarrow=\{\mathbf w:\mathbf w I_d\mathbf x\text{ がすべての }\mathbf x\in C\text{ について成り立つ}\}
\]
と置く。
閉対象集合 $A^{\uparrow\downarrow}$、同値に formal concept $(A^{\uparrow\downarrow},A^\uparrow)$ 全体が Wurm の $\SCL_d(L)$ をなす。
タプルの成分ごとの連接を行い、その後 closure を取ることでモノイド演算が得られ、さらに residual により residuated lattice となる~\cite{WurmSCLTuples}。

タプル $\mathbf w$ に対し
\[
\gamma_d(\mathbf w)
:=(\{\mathbf w\}^{\uparrow\downarrow},\{\mathbf w\}^\uparrow)
\]
をその\emph{対象概念}と書き、
\[
\Obj_d(L):=\{\gamma_d(\mathbf w):\mathbf w\in(\Sigma^*)^d\}
\]
を主層／対象層とする。
$d=1$ では Clark の結果~\cite[Proposition~3.19]{ClarkSCL} により
\[
\gamma_1(u)=\gamma_1(v)
\quad\Longleftrightarrow\quad
u\equiv_L v
\]
である。
また singleton concept の積は
$\gamma_1(u)\circ\gamma_1(v)=\gamma_1(uv)$
なので、次の標準的なモノイド同型を得る。

\begin{proposition}[単項主モノイド]
\label{scl:prop:unary-principal}
任意の言語 $L$ について、単項対象概念のモノイドは標準的に
\[
\Obj_1(L)\cong \Sigma^*/{\equiv_L}
\]
である。
$L$ が正則ならば、これは有限統語モノイド $T_L$ である。
\end{proposition}

\subsection{主 intent の重なり}

タプル $\vec x\in(\Sigma^*)^d$ に対し
\[
\D_L^{(d)}(\vec x)
:=\{(u_0,\ldots,u_d):u_0x_1u_1\cdots x_du_d\in L\}
\]
と書く。
これは主概念 $\gamma_d(\vec x)$ の intent そのものである。
$d=1$ では $\D_L(x)$ と書く。

\begin{definition}[主重なりグラフ]
$d\ge1$ に対し、$\Om_d(L)$ を頂点集合 $\Obj_d(L)$ をもつ単純グラフとする。
異なる対象概念 $A,B$ が隣接するのは、かつそのときに限り
\[
\operatorname{Int}(A)\cap\operatorname{Int}(B)\ne\varnothing
\]
とする。
\end{definition}

したがって隣接とは、分布的には異なる二つのタプル概念が、少なくとも一つの共通受理文脈によって許されることを意味する。
有限モノイドへの準同型 $h:\Sigma^*\to M$ は、タプル上に成分ごとの色
\[
h^{(d)}(w_1,\ldots,w_d)
:=(h(w_1),\ldots,h(w_d))\in M^d
\]
を誘導する。
安全性条件は、主重なり衝突を潰すことなく、これらの成分ごとの色のうちどれを同一視できるかを問う。

\begin{definition}[アリティ $f$ までの有限モノイド安全性]
\label{scl:def:safety}
$h:\Sigma^*\to M$ を有限モノイドへの準同型とし、$f\ge1$ とする。
すべての $1\le d\le f$ に対して、$\Om_d(L)$ の相異なる隣接頂点を生成する二つのタプルが同じ成分ごとの $h$-型を受け取らないとき、
\[
\operatorname{Safe}_f(h;L)
\]
と書く。
同値に、相異なる隣接主概念を二つ取ったとき、それぞれから生成タプルをどのように一つずつ選んでも、両者の成分ごとの $h$-型が一致してはならない。
同義語として、このとき $L$ は $(f,h)$-tuple-substitutable であるともいう。
これは別々の二概念ではない。
本稿では同じ安全性述語に対する二つの名称であり、前者は SCL 幾何を、後者はガード付き代入タスクを強調する。
\end{definition}

$\gamma_d(\vec x)$ の intent はちょうど $\D_L^{(d)}(\vec x)$ なので、$\Om_d(L)$ の辺とは、まさに「重なりはあるが同値ではない」組である。
すなわち、二つのタプルは少なくとも一つの受理文脈を共有するが、別の文脈によって区別される。
したがって $\operatorname{Safe}_f(h;L)$ は操作的にはすでに次のように読める。
有限な $h$-型が等しいことによってこのようなタプルを同一視してよいのは、それらの完全なタプル分布が一致するときに限る。
正則言語では、同じ witness を合同形式および有限関係形式でも表せる。
本稿では表現面ごとの用語を維持し、圧縮写像には\emph{safe}、合同および関係的射には\emph{separating}という語を用いる。
それらの厳密な同値性は Safety Interface Theorem、すなわち定理~\ref{alg:thm:safety-interface} に記す。

\begin{remark}[なぜ主層が代入層なのか]
\label{scl:rem:principal-canonical}
主概念への制限は、SCL を恣意的に切り詰めたものではない。
タプル文脈へ実際に挿入される対象は語および語のタプルであり、Clark--Wurm 束におけるそれらの像が、まさに主概念 $\gamma_d(\vec x)$ である。
一般の閉概念は、追加の合成入力を表すのではなく、タプルの Galois 閉包された集合体を表す。
したがって定義~\ref{scl:def:safety} は、ガード付きタプル代入によって強制される正確な衝突幾何を取り出している。
本稿における「標準性」の主張はこのタスクに相対的なものであり、主重なり圧縮が概念束全体に対する唯一の意味ある圧縮概念だと主張するものではない。
Table~\ref{tab:scl-dictionary} に対応関係をまとめる。
\end{remark}

\begin{table}[t]
\centering
\caption{主 SCL 幾何とタプル代入言語の対応。}
\label{tab:scl-dictionary}
\begin{tabular}{@{}ll@{}}
\toprule
\textbf{SCL 側の表現} & \textbf{語／文脈側の表現}\\
\midrule
主対象概念 $\gamma_d(\vec x)$ & $\vec x$ のタプル分布類\\
$\gamma_d(\vec x)$ の intent & $\D_L^{(d)}(\vec x)$\\
異なる主概念 & 異なるタプル分布\\
intent の重なり & 共通の受理タプル文脈\\
$\Om_d(L)$ の辺 & unsafe なタプル対\\
有限モノイド圧縮 & 座標ごとの衝突分離\\
\bottomrule
\end{tabular}
\end{table}

\subsection{圧縮数}

\[
\cmp_f(L)
:=\min\{|\operatorname{im}(h)|:h\text{ はアリティ }f\text{ まで safe}\}
\]
と定め、有限 witness が存在しない場合は値を $\infty$ とする。
これを、アリティ $f$ までの主 SCL 幾何の\emph{有限モノイド圧縮数}と呼ぶ。

合成的両立性を完全に無視した比較量も有用である。

\begin{remark}[companion MCFG の観測パラメータとの関係]
\label{rem:obs-cmp-reindex}
companion MCFG manuscript~\cite{KuriyamaMCFG2026} では、$L$ を orientation-aware な $(f,h)$-tuple-substitutable にする有限観測の像の最小サイズを $\obs_f(L)$ と書く。
そこでは名前付き文脈が置換 sector に分割されるのに対し、本稿で用いる Clark--Wurm 規約ではタプル座標の表面順序を固定する。

companion manuscript の条件は sector ごとの条件である。
すなわち substitutability は、異なる sector に属する文脈同士を比較するのではなく、各固定置換 sector 内で検査される。
任意の $\sigma\in S_d$ に対して、添字の付け替え
\[
(x_1,\ldots,x_d)
\longmapsto
(x_{\sigma(1)},\ldots,x_{\sigma(d)})
\]
は $\sigma$-sector と本稿の固定順序文脈とを同一視する。
固定順序での安全性条件はすべてのタプル上で量化されるため、このように添字を付け替えた任意の対にも同様に適用される。
さらに、二つのタプルに同じ置換を施しても、成分ごとの $h$-値の等しさは保存される。

したがって安全性条件は sector ごとに一致し、同じ言語に二つの記法を適用する限り
\[
\obs_f(L)=\cmp_f(L)
\]
である。
本稿ではこの数を学習側の情報パラメータではなく SCL 圧縮不変量として研究するため、記号 $\cmp_f$ を用いる。
\end{remark}

\begin{proposition}[幾何学的下界]
\label{scl:prop:chromatic-lower}
任意の言語 $L$ と任意の $d\ge1$ に対して
\[
\left\lceil\chi(\Om_d(L))^{1/d}\right\rceil\le \cmp_d(L)
\]
が成り立つ。
\end{proposition}
\begin{proof}
像の値を $m$ 個もつ witness $h$ を取る。
各対象概念 $A\in\Obj_d(L)$ に対し、$A$ を生成するあるタプルが実現する成分ごとの $h$-値を一つ割り当てる。
この選択は標準的である必要はない。
それでも、隣接する二概念が $M^d$ の同じベクトルを受け取ったなら、選ばれた代表タプルはそれぞれの概念と同じ主 intent をもつ。
したがって両者のタプル分布は異なる一方で重なりをもち、すべての座標で潰される unsafe pair を形成することになり、安全性に反する。
よってこれは高々 $m^d$ 色を用いる proper coloring である。
したがって $\chi(\Om_d(L))\le m^d$。
$m$ は正整数なので
$\lceil\chi(\Om_d(L))^{1/d}\rceil\le m$ である。
witness 全体で最小化すれば主張を得る。
同じ議論から、$\chi(\Om_d(L))$ が無限なら有限 witness は存在せず、したがって $\cmp_d(L)=\infty$ であることも従う。
\end{proof}

\subsection{正則言語に対する有限 SCL 還元}

$L$ が正則であるとし、$\eta:\Sigma^*\twoheadrightarrow T$ をその統語射、$P=\eta(L)$ とする。
$k\ge1$ に対し、成分ごとの拡張を
$\eta^{(k)}:(\Sigma^*)^k\to T^k$
と書く。
有限 formal context
\[
\mathcal K_d(T,P)=\bigl(T^d,T^{d+1},I_{d,T,P}\bigr)
\]
を、
\[
(t_1,\ldots,t_d)\ I_{d,T,P}\ (q_0,\ldots,q_d)
\quad\Longleftrightarrow\quad
q_0t_1q_1\cdots t_dq_d\in P
\]
によって定める。
対象上の $\eta^{(d)}$ と文脈上の $\eta^{(d+1)}$ はともに全射であり、incidence relation を両方向に保存する。
したがって失われる情報は、重複した行と列だけである。

\begin{proposition}[有限概念への還元]
\label{scl:prop:finite-reduction}
任意の正則言語 $L$ と有限 $d$ に対し、$\SCL_d(L)$ は $\mathcal K_d(T_L,P)$ の概念束と同型である。
特に、$\mathbf t=(t_1,\ldots,t_d)\in T^d$ の主 intent は
\[
\Phi_d(\mathbf t)
=\{(q_0,\ldots,q_d)\in T^{d+1}:q_0t_1q_1\cdots t_dq_d\in P\}
\]
である。
\end{proposition}
\begin{proof}
具体的なタプル $\mathbf w$ と文脈 $\mathbf x$ に対し、incidence は $\mathcal K_d(T,P)$ における $\eta^{(d)}(\mathbf w)$ と $\eta^{(d+1)}(\mathbf x)$ の incidence と同値である。
両写像は全射である。
有限 context の Galois 演算子を $\uparrow_T,\downarrow_T$ と書くと、任意の対象集合 $A\subseteq(\Sigma^*)^d$ と文脈集合 $C\subseteq(\Sigma^*)^{d+1}$ に対し
\[
A^\uparrow
=(\eta^{(d+1)})^{-1}\!\bigl((\eta^{(d)}(A))^{\uparrow_T}\bigr),
\qquad
C^\downarrow
=(\eta^{(d)})^{-1}\!\bigl((\eta^{(d+1)}(C))^{\downarrow_T}\bigr)
\]
が成り立つ。
したがって無限 context の任意の閉 extent と intent は、対応する統語写像に関して saturated であり、有限 context における一意な閉 extent または intent の逆像となる。
逆向きは全射性から従う。
ゆえに二つの概念束は同型である。
\end{proof}

以下の各節では、この SCL 定式化を合同および関係的射の最適化へ変換し、その後アリティ挙動を分類する。

\section{語レベルの衝突と合成的圧縮}

\label{sec:conflict}

本稿に現れるモノイドはすべて単位元をもち、すべての準同型は単位元を保存するものとする。
有限アルファベット $\Sigma$ を固定する。
Section~\ref{sec:scl-geometry} で導入した固定順序 Clark--Wurm タプル文脈と充填記法 $E[\vec x]$ を引き続き用いる。
したがってタプル分布 $\D_L^{(d)}(\vec x)$ は、同値に
$E[\vec x]\in L$
を満たすタプル文脈 $E$ 全体の集合である。

語レベルで新たな安全性述語を導入するわけではない。
定義~\ref{scl:def:safety} の同義表現「$L$ は $(f,h)$-tuple-substitutable である」を用いる。
本節の目的は、この条件から得られる意味論的な衝突節を明示することである。
後に定理~\ref{alg:thm:safety-interface} において、同じ述語が有限統語形式および関係形式と一致することを示す。

\begin{proposition}[認識から圧縮への緩和]
\label{conf:prop:recognition-relaxation}
$h:\Sigma^*\to M$ を $L$ を認識するモノイド射とし、ある $P\subseteq M$ に対して
$L=h^{-1}(P)$
であるとする。
このとき $h$ は任意の有限 $f$ に対して $(f,h)$-tuple substitutability の witness である。
したがって $L$ が正則ならば
\[
\cmp_f(L)\le |T_L|
\qquad(f\ge1)
\]
が成り立つ。ただし $T_L$ は $L$ の統語モノイドである。
\end{proposition}

\begin{proof}
ある有限 $d$ について
$h^{(d)}(\vec x)=h^{(d)}(\vec y)$
とする。
任意のタプル文脈
$E=u_0\blank_1u_1\cdots\blank_du_d$
に対し、$h$ が連接と両立することから
\[
h(E[\vec x])=h(E[\vec y])
\]
である。
$L$ への所属は $h$-値だけで決まるので、
$E[\vec x]\in L$
であることと
$E[\vec y]\in L$
であることは同値である。
したがって二つのタプル分布は等しく、$h$ は任意の有限アリティで safe である。
$L$ が正則ならば、この主張を統語射
$\eta_L:\Sigma^*\twoheadrightarrow T_L$
に適用すればよい。
\end{proof}

\begin{remark}[厳密認識と十分な圧縮]
命題~\ref{conf:prop:recognition-relaxation} により、緩和の意味は文字どおり明確になる。
厳密な recognizer は、圧縮問題に対する制限付きの実行可能解のクラスをなす。
したがって不変量 $\cmp_f$ は、タスクを厳密認識からアリティ $f$ までの文脈代入安全性へ弱めたとき、有限の合成的情報をどれだけ捨てられるかを測る。
以下の例では、すべての有限アリティを要求してもこの緩和による差が任意に大きくなりうる一方、別の族では階層がすでにアリティ2で安定することを示す。
中心をもたない有限群の核言語では、二項レベルだけで統語モノイド全体の大きささえ回復する。
\end{remark}

\begin{definition}[文脈衝突／unsafe タプル対]
$d\ge1$ に対し
\[
\mathfrak U_d(L)
:=\left\{(\vec x,\vec y):
\D_L^{(d)}(\vec x)\cap\D_L^{(d)}(\vec y)\neq\varnothing,
\ \D_L^{(d)}(\vec x)\neq\D_L^{(d)}(\vec y)
\right\}
\]
と置く。
その元を $L$ の\emph{文脈衝突}、または\emph{アリティ $d$ の unsafe タプル対}と呼ぶ。
\end{definition}

この定義の二条件は異なる役割をもつ。
完全な分布が等しければ、代入は意味論的に安全であり、分離の必要はない。
分布が互いに素ならば、二つのタプルが同時に許される成功文脈が存在しないため、ガード付き代入原理はそもそも両者を比較しない。
したがって unsafe pair は、解消しなければならない「重なりはあるが同値ではない」配置とちょうど一致する。
同値に、$h$ が $\cmp_f$ の witness であることは、すべての
$(\vec x,\vec y)\in\mathfrak U_d(L)$, $d\le f$ に対して
\[
h(x_1)\neq h(y_1)
\ \vee\ \cdots\ \vee\
h(x_d)\neq h(y_d)
\]
が成り立つことと同値である。

\subsection{アリティ1の衝突グラフ}

\begin{definition}[文脈衝突グラフ]
言語 $L\subseteq\Sigma^*$ に対し、$\Gamma_L$ を頂点集合 $\Sigma^*$ をもつ単純グラフとする。
異なる語 $x,y$ が辺 $\{x,y\}$ をなすのは、それらが unsafe な単項対であるとき、同値に
\[
\D_L(x)\cap\D_L(y)\neq\varnothing
\quad\text{かつ}\quad
\D_L(x)\neq\D_L(y)
\]
のときとする。
\end{definition}

\begin{theorem}[合成的彩色による特徴づけ]
\label{conf:thm:coloring}
任意の言語 $L$ に対して
\[
\cmp_1(L)
=
\min\left\{
|\operatorname{im}(h)|:
\begin{array}{l}
h:\Sigma^*\to M\text{ は有限モノイドへの準同型であり、}\\
\Gamma_L\text{ の proper coloring でもある}
\end{array}
\right\}.
\]
特に
\[
\chi(\Gamma_L)\le\cmp_1(L).
\]
\end{theorem}

\begin{proof}
アリティ1では、準同型が tuple substitutability の witness であることは、unsafe pair のどれも同じ値へ送られないこととちょうど同値である。
これは $\Gamma_L$ の proper-coloring 条件そのものである。
準同型という制約を忘れれば通常の proper coloring が残るので、不等式も従う。
\end{proof}

したがって通常の彩色数は衝突幾何だけを測るのに対し、$\cmp_1$ はその幾何を連接と両立するように彩色するコストを測る。
両者の差は、いわば合成可能性の代価である。

\begin{remark}[合成可能性に対する真の追加コスト]
\label{conf:rem:strict-price}
$L=\{a,a^3\}\subseteq\{a\}^*$ とする。
非自明な単項衝突辺は
$\{\varepsilon,a^2\}$ と $\{a,a^3\}$
だけなので、$\chi(\Gamma_L)=2$ である。
しかし $\cmp_1(L)=3$ となる。

像が1元しかなければ両方の衝突辺を潰すため witness にはなれない。
したがって像がちょうど2値の witness が存在しないことを示せばよい。
そのような準同型は $x=h(a)$ によって決まり、必ず $x\ne1$ である。
$x$ が生成する2元モノイドには、群の場合 $x^2=1$ と、冪等の場合 $x^2=x$ の二通りしかない。
前者では $\varepsilon$ と $a^2$ が同一視され、後者では $x^3=x$ なので $a$ と $a^3$ が同一視される。
よって両方の衝突辺を分離する2元圧縮写像は存在しない。
一方、$a$ を巡回群 $C_3$ の生成元へ送れば両辺を分離できる。
したがって
\[
\chi(\Gamma_L)=2<3=\cmp_1(L).
\]
これは通常のグラフ彩色だけでは捉えられない最小型の現象であり、追加コストは彩色に合成可能性を要求することだけから生じる。
\end{remark}

\subsection{正則言語：彩色から統語意味論へ}

正則言語 $L$ に対し、$\eta_L:\Sigma^*\twoheadrightarrow T_L$ を統語射とする。
圧縮写像は $L$ を認識する必要はなく、統語類上で一定である必要もない。
要求されるのは文脈衝突を分離することだけである。

単項文脈分布は統語類上で一定なので、無限グラフ $\Gamma_L$ は有限な意味論的核をもつ。
$T_L$ 上の\emph{統語衝突グラフ} $\Gamma_L^{\mathrm{syn}}$ を次のように定義する。
異なる元 $s,t$ を、$\eta_L(x)=s$, $\eta_L(y)=t$ を満たす任意の代表元 $x,y$ が unsafe な単項対をなすとき、かつそのときに限り辺で結ぶ。

\begin{proposition}[正則言語の有限衝突核]
\label{conf:prop:syntactic-conflict-graph}
正則言語 $L$ に対し、$\Gamma_L$ は $\Gamma_L^{\mathrm{syn}}$ の各統語頂点を語代表元からなる独立集合に置き換えることで得られる。
特に
\[
\chi(\Gamma_L)=\chi(\Gamma_L^{\mathrm{syn}}).
\]
\end{proposition}

\begin{proof}
同じ統語類に属する語は同一の両側文脈分布をもつため、互いに隣接しない。
異なる二つの統語類に対しては、文脈分布の交わりが非空かどうか、および分布が異なるかどうかは類だけに依存する。
したがって、二類の代表元のすべての組が隣接するか、どの組も隣接しないかのいずれかである。
ゆえに $\Gamma_L$ は有限統語衝突グラフの blow-up であり、各頂点を非空独立集合へ blow-up しても彩色数は変わらない。
\end{proof}

\begin{corollary}[単項 SCL の重なりと統語衝突の一致]
\label{conf:cor:omega-syntactic}
命題~\ref{scl:prop:unary-principal} の標準同一視
$\Obj_1(L)\cong T_L$
のもとで、正則言語 $L$ について
\[
\Om_1(L)=\Gamma_L^{\mathrm{syn}}
\]
である。
したがって
\[
\chi(\Gamma_L)
=\chi(\Gamma_L^{\mathrm{syn}})
=\chi(\Om_1(L)).
\]
\end{corollary}
\begin{proof}
語 $x$ によって表される単項主概念の intent は $\D_L(x)$ である。
定義~\ref{scl:def:safety} により、$\Om_1(L)$ の隣接性は対応する統語類間の unsafe 関係と正確に一致する。
残る彩色数の等式は命題~\ref{conf:prop:syntactic-conflict-graph} である。
\end{proof}

\begin{definition}[商制限付き圧縮サイズ]
正則言語 $L$ に対し、$\qcmp_f(L)$ を、$T_L$ の商モノイド $Q$ のうち、
\[
\pi\circ\eta_L:\Sigma^*\to Q
\]
が $(f,\pi\circ\eta_L)$-tuple substitutability の witness となるものの最小サイズとする。
\end{definition}

\begin{proposition}[衝突--統語複雑性の梯子]
\label{conf:prop:ladder}
任意の正則言語 $L$ に対して
\[
\chi(\Gamma_L)
\le \cmp_1(L)
\le \qcmp_1(L)
\le |T_L|.
\]
\end{proposition}

\begin{proof}
最初の不等式は定理~\ref{conf:thm:coloring} である。
第二の不等式は、商圧縮が有限圧縮写像の制限されたクラスだから従う。
第三の不等式については、$T_L$ の恒等商を取ればよい。
二つの語が同じ統語値をもつなら、任意の両側文脈内で代入しても所属は保存されるので、統語射自身が witness である。
\end{proof}

この梯子は四つの資源を分離している。
すなわち、無制約な衝突彩色、合成的な衝突彩色、統語意味論の商に基づく抽象化、そして完全な統語モノイドである。
四つの不等式はいずれも真に厳しくなりうるし、中央のギャップは任意に大きくなりうる。

\begin{proposition}[関係的圧縮と商圧縮の差が非有界となる真の梯子]
\label{conf:prop:strict-ladder-family}
$n\ge2$ に対して
\[
K_n=\{a^n,a^{n+2}\}\subseteq\{a\}^*
\]
とする。
このとき任意の $f\ge1$ について
\[
\chi(\Gamma_{K_n})=2
<3=\cmp_f(K_n)
<n+2=\qcmp_f(K_n)
<n+4=|T_{K_n}|.
\]
特に
$\qcmp_f(K_n)-\cmp_f(K_n)=n-1$
であり、有限圧縮を統語モノイドの商へ制限することには非有界な加法的コストがある。
\end{proposition}

\begin{proof}
$0\le i\le n+2$ に対し $a^i$ の統語類を $t_i$ と書き、$j\ge n+3$ のすべての $a^j$ が属する共通類を $\bot$ と書く。
これらがちょうど統語類の全体である。
長さ $n+2$ 以下の異なる二語は、一方を長さ $n$ または $n+2$ に完成させる文脈によって区別できる一方、それより長い語はすべて空の分布をもつからである。
したがって $|T_{K_n}|=n+4$ であり、積は truncated addition、すなわち $i+j\le n+2$ なら
$t_it_j=t_{i+j}$、それ以外では積は $\bot$ である。

非自明な単項衝突辺は正確に
\[
\{t_i,t_{i+2}\},
\qquad 0\le i\le n
\]
である。
したがって有限衝突核は、偶数添字上の path と奇数添字上の path の直和であり、
$\chi(\Gamma_{K_n})=2$ となる。

$a$ を巡回群 $C_3$ の生成元 $g$ へ送れば、$g^i\ne g^{i+2}$ なので、すべての上記辺を分離する。
よって $\cmp_1(K_n)\le3$。

衝突グラフには辺があるので、像が1元の写像は unary-safe ではない。
像が2元のものも使えない。
$a$ の像を $x$ とし、$x$ が生成する像が2元だとすると、$x^2=1$ または $x^2=x$ のいずれかである。
前者ではすべての対 $a^i,a^{i+2}$ が同一視され、後者では $a$ と $a^3$ が同一視される。
したがって $\cmp_1(K_n)=3$ である。
$T_{K_n}$ は可換なので、後に証明する可換崩壊定理、定理~\ref{group:thm:commutative-collapse} により、任意の $f\ge1$ について
$\cmp_f(K_n)=3$ である。

残るのは商制限値の計算である。
unary-safe な任意の商
$\pi:T_{K_n}\twoheadrightarrow Q$
を取り、$y=\pi(t_1)$, $z=\pi(\bot)$ とする。
$y^r=z$ となる最小の $r$ を取る。
$y^{n+3}=z$ なので、このような $r$ は存在する。
もし $r\le n$ なら
$\pi(t_r)=z=\pi(t_{r+2})$
となり、衝突辺が潰れてしまう。
したがって $r\ge n+1$。
さらに
$1,y,\ldots,y^{r-1},z$
は相異なる。
実際、$0\le i<j<r$ で $y^i=y^j$ なら、両辺に $y^{r-j}$ を掛けて
$y^{i+r-j}=z$
を得るが、指数 $i+r-j<r$ となり $r$ の最小性に反する。
よって
$|Q|\ge r+1\ge n+2$。

逆に、tail
$\{t_{n+1},t_{n+2},\bot\}$
を一つに同一視し、$t_0,\ldots,t_n$ を互いに区別して残す。
これはサイズ $n+2$ の商モノイドであり、どの衝突辺も潰さない。
したがって
$\qcmp_1(K_n)=n+2$。
さらに可換崩壊により同じ商は任意の有限アリティまで safe なので、すべての $f\ge1$ に対し
$\qcmp_f(K_n)=n+2$ である。
\end{proof}

\begin{remark}[なぜ最適圧縮は統語上で本質的に関係的なのか]
\label{conf:rem:genuinely-relational}
上の最適写像 $h(a)=g\in C_3$ に対して、統語モノイド上の標準関係は
\[
\rho_h(t_i)=\{g^i\}\quad(0\le i\le n+2),
\qquad
\rho_h(\bot)=C_3
\]
を満たす。
したがって一つの統語元が三つの圧縮値をもちうる。
有限圧縮は dead な統語類の内部に剰余情報を保持しており、その情報は $T_{K_n}$ の任意の関数的商では失われる。
これは抽象的な selector theorem に依存せず、商意味論ではなく関係意味論が正しい有限インターフェースであることの具体的理由を与える。
\end{remark}

Section~\ref{sec:graph-universality} のグラフ実現定理により、最初の二段階が任意の有限衝突幾何上で一致する場合もあることを後に示す。
以下の構造理論にとって重要な次の段階は、高次アリティである。

\subsection{合成的分離系としての高次アリティ}

$d>1$ では unsafe pair はもはや語の色の間の単一グラフ辺ではない。
それは選言節
\[
h(x_1)\neq h(y_1)
\ \vee\ \cdots\ \vee\
h(x_d)\neq h(y_d)
\]
を課す。
したがって $\cmp_f(L)$ は、$\Sigma^*$ の標準的彩色が、言語から誘導される幅 $f$ 以下のすべての分離節を満たすような有限モノイドの最小サイズである。
グラフの場合は、この制約系の幅1の影にすぎない。
次節の有限指数合同定理により、この定式化を代数的に厳密化する。

\section{衝突解消の代数的特徴づけ}
\label{sec:algebraic-conflict}

Section~\ref{sec:conflict} の衝突節は、厳密な代数形式をもつ。
自由モノイド上では、有限圧縮写像とは、すべての衝突を合成的に彩色する有限指数合同と同じものである。
正則言語では、この最適化を統語モノイドへ押し下げられるが、一般には関数的な商ではなく関係的射を経由する必要がある。

自由モノイド上の合同定式化は初等的だが有用である。
それは最適化問題を内在的な代数形式で記録する。
正則言語における本質的な段階は、その後に自由モノイドから有限の点付き統語モノイド上の関係的射へ移行することである。

\subsection{自由モノイド上の分離節}

最初の代数的インターフェースは、圧縮写像の核合同である。
ここでは実行可能性の概念だけを導入する。
最適化定理は Safety Interface Theorem の後で述べ、共通インターフェースから直ちに従う形にする。

\begin{definition}[分離合同]
有限指数の両側モノイド合同
\[
\theta\subseteq\Sigma^*\times\Sigma^*
\]
が $L$ に対して\emph{$f$-分離的}であるとは、任意の
\[
(\vec x,\vec y)\in\mathfrak U_d(L),
\qquad
1\le d\le f
\]
に対し、ある座標 $i\in\{1,\ldots,d\}$ が存在して
$x_i\not\equiv_\theta y_i$
となることをいう。
\end{definition}

\subsection{正則言語に対する有限統語 profile}

以下、$L$ は正則であると仮定する。
$\eta:\Sigma^*\twoheadrightarrow T$ をその統語射とし、
$P:=\eta(L)\subseteq T$
とする。したがって $L=\eta^{-1}(P)$ である。
古典的な pointed-monoid（または $P$-monoid）の枠組み~\cite{Sakarovitch1979,Pin1981} に従い、$(T,P)$ を $L$ の\emph{点付き統語モノイド}と呼ぶ。

命題~\ref{scl:prop:finite-reduction} で導入した有限タプル文脈 profile
\[
\Phi_d(\mathbf t)
=\left\{(q_0,\ldots,q_d)\in T^{d+1}:q_0t_1q_1\cdots t_dq_d\in P\right\}
\qquad(\mathbf t\in T^d)
\]
を思い出す。

\[
E=u_0\blank_1u_1\cdots\blank_du_d
\]
ならば
\[
\eta_{\mathrm{ctx}}^{(d)}(E)
:=(\eta(u_0),\ldots,\eta(u_d))\in T^{d+1}
\]
と書く。
この記法はタプル文脈の有限統語 profile のためだけに用いる。

\begin{lemma}[有限 profile の移送]
\label{alg:lem:profile}
$\vec x,\vec y\in(\Sigma^*)^d$ とし、
\[
\mathbf s:=\eta^{(d)}(\vec x),
\qquad
\mathbf t:=\eta^{(d)}(\vec y)
\]
と置く。このとき次が成り立つ。
\begin{enumerate}[label=(\roman*),leftmargin=*]
\item 任意のタプル文脈 $E$ に対して
\[
E\in\D_L^{(d)}(\vec x)
\iff
\eta_{\mathrm{ctx}}^{(d)}(E)\in\Phi_d(\mathbf s).
\]
\item
\[
\D_L^{(d)}(\vec x)=\D_L^{(d)}(\vec y)
\iff
\Phi_d(\mathbf s)=\Phi_d(\mathbf t).
\]
\item
\[
\D_L^{(d)}(\vec x)\cap\D_L^{(d)}(\vec y)\neq\varnothing
\iff
\Phi_d(\mathbf s)\cap\Phi_d(\mathbf t)\neq\varnothing.
\]
\end{enumerate}
\end{lemma}

\begin{proof}
(i) について、
$E=u_0\blank_1u_1\cdots\blank_du_d$
とすると
\[
\eta(E[\vec x])
=\eta(u_0)s_1\eta(u_1)\cdots s_d\eta(u_d).
\]
$L=\eta^{-1}(P)$ なので、この語が $L$ に属することは
$\eta_{\mathrm{ctx}}^{(d)}(E)\in\Phi_d(\mathbf s)$
とちょうど同値である。

(ii) の順方向については、(i) に加えて、任意の抽象 profile
$(q_0,\ldots,q_d)\in T^{d+1}$
が具体的な文脈代表元をもつことを用いる。
$\eta$ は全射なので、$\eta(u_j)=q_j$ を満たす $u_j\in\Sigma^*$ を選び、
$u_0\blank_1u_1\cdots\blank_du_d$
を作ればよい。
したがって $\Phi_d(\mathbf s)$ と $\Phi_d(\mathbf t)$ の一方だけに属する profile があれば、それに対応する具体的文脈は二つのタプル分布の一方だけに属する。
逆方向は (i) から直ちに従う。

(iii) について、共通の具体的受理文脈があれば (i) により共通 profile が得られる。
逆に抽象 profile が両方の有限 profile 集合に属するなら、上と同様に各 $u_j$ を選ぶことで、その一つの名前付き文脈が (i) により両方のタプルを受理する。
\end{proof}

\begin{definition}[unsafe 統語タプル対]
$1\le d\le f$ に対して
\[
\mathcal U_d(T,P)
:=
\left\{
(\mathbf s,\mathbf t)\in T^d\times T^d:
\begin{array}{l}
\Phi_d(\mathbf s)\cap\Phi_d(\mathbf t)\neq\varnothing,\\
\Phi_d(\mathbf s)\neq\Phi_d(\mathbf t)
\end{array}
\right\}
\]
と定める。
\end{definition}

補題~\ref{alg:lem:profile} により、これは意味論的 unsafe-pair 関係の有限統語像と正確に一致する。
以下、語タプルの意味論的 unsafe pair には $\mathfrak U_d(L)$ を、その有限統語 profile 版には $\mathcal U_d(T,P)$ を用いる。

\subsection{圧縮写像が誘導する標準関係的射}

\begin{definition}[モノイド関係的射]
$T,M$ をモノイドとする。
本稿では単位元保存のモノイド規約を用いる。
\emph{モノイド関係的射}
$\rho:T\relto M$
とは、各 $t\in T$ に非空集合 $\rho(t)\subseteq M$ を対応させ、
\[
1_M\in\rho(1_T),
\qquad
\rho(s)\rho(t)\subseteq\rho(st)
\quad(s,t\in T)
\]
を満たすものをいう。
同値に、そのグラフ
\[
\operatorname{Gr}(\rho)
:=\{(t,m):m\in\rho(t)\}
\]
は $T\times M$ の部分モノイドであり、第一射影が全射である。
第二射影が全射であることは要求しない。
使われていない終域値は常に削除でき、定理~\ref{alg:thm:relational-characterization} の証明で明示的にそうする。
\end{definition}

\begin{definition}[圧縮写像が誘導する標準関係的射]
$h:\Sigma^*\to M$ を有限圧縮写像とする。
$\rho_h:T\relto M$ を
\[
\rho_h(t):=\{h(w):\eta(w)=t\}
\]
で定める。
\end{definition}

\begin{lemma}[標準関係は関係的射である]
\label{alg:lem:canonical-rel}
$\rho_h:T\relto M$ はモノイド関係的射である。
そのグラフは正確に
\[
\operatorname{im}(\eta,h)
=
\{(\eta(w),h(w)):w\in\Sigma^*\}
\subseteq T\times M
\]
である。
\end{lemma}

\begin{proof}
全射な統語射のもとで各 $t\in T$ は原像をもつので、$\rho_h(t)\neq\varnothing$ である。
また
$1_M=h(\varepsilon)\in\rho_h(1_T)$。
$m\in\rho_h(s)$, $n\in\rho_h(t)$ とする。
\[
\eta(u)=s,\ h(u)=m,
\qquad
\eta(v)=t,\ h(v)=n
\]
を満たす語 $u,v$ を選ぶ。
すると
$\eta(uv)=st$, $h(uv)=mn$
なので $mn\in\rho_h(st)$ である。
グラフに関する等式は定義から直ちに従う。
\end{proof}

\begin{definition}[unsafe を分離する関係的射]
\label{alg:def:rel-separating}
関係的射 $\rho:T\relto M$ が $(T,P)$ に対して\emph{$f$-分離的}であるとは、任意の
\[
(\mathbf s,\mathbf t)\in\mathcal U_d(T,P),
\qquad 1\le d\le f
\]
に対して、ある座標 $i$ が存在し
\[
\rho(s_i)\cap\rho(t_i)=\varnothing
\]
となることをいう。
\end{definition}

\begin{definition}[safe compatible tolerance]
\label{alg:def:safe-tolerance}
$(T,P)$ を有限点付きモノイドとする。
関係 $\tau\subseteq T\times T$ が\emph{compatible tolerance} であるとは、反射的かつ対称的で、さらに
\[
x\tau y,
\quad
u\tau v
\Longrightarrow
xu\tau yv
\]
を満たすことをいう。
任意の
$(\mathbf s,\mathbf t)\in\mathcal U_d(T,P)$, $d\le f$
がある座標で分離される、すなわちある $i$ に対し
$s_i\not\tau t_i$
となるとき、$\tau$ は $(T,P)$ に対して\emph{$f$-safe} であるという。
\end{definition}

\begin{lemma}[関係的重なりインターフェース]
\label{alg:lem:relational-overlap}
任意の関係的射 $\rho:T\relto M$ は compatible tolerance
\[
x\tau_\rho y
\quad\Longleftrightarrow\quad
\rho(x)\cap\rho(y)\ne\varnothing
\tag{T}
\]
を誘導する。
さらに
\[
\rho\text{ が }(T,P)\text{ に対して }f\text{-分離的}
\quad\Longleftrightarrow\quad
\tau_\rho\text{ が }(T,P)\text{ に対して }f\text{-safe}
\tag{RT}
\]
である。
\end{lemma}

\begin{proof}
反射性と対称性は $\rho$ の各 fiber が非空であることから従う。
$x\tau_\rho y$ かつ $u\tau_\rho v$ とする。
\[
m\in\rho(x)\cap\rho(y),
\qquad
n\in\rho(u)\cap\rho(v)
\]
を選ぶ。
すると
$mn\in\rho(xu)\cap\rho(yv)$
なので compatibility が従う。
同値性 \textup{(RT)} は二つの定義から直ちに従う。
unsafe tuple pair が $\rho$ によって分離されないことは、そのすべての座標対が $\tau_\rho$ で関係づけられていることと正確に同値である。
\end{proof}

\begin{theorem}[Safety Interface Theorem]
\label{alg:thm:safety-interface}
$L\subseteq\Sigma^*$ を任意の言語とし、$h:\Sigma^*\to M$ を有限モノイドへの準同型とする。
任意の $f\ge1$ に対して、以下の \textup{(i)}--\textup{(iv)} は同値である。
さらに $L$ が正則で、$\eta:\Sigma^*\twoheadrightarrow T$ を統語射、$P=\eta(L)$ と書くとき、これらは \textup{(v)} とも同値である。
\begin{enumerate}[label=(\roman*),leftmargin=*]
\item \textbf{主 SCL インターフェース：}
$\operatorname{Safe}_f(h;L)$ が成り立つ。
同値に、任意の $d\le f$ について、$\Om_d(L)$ の相異なる隣接主概念を生成する二つのタプルが、同じ成分ごとの $h$-型をもつことはない。
主概念から $M^d$ への写像の存在を主張したり要求したりしているわけではない。

\item \textbf{ガード付き語／文脈インターフェース：}
任意の $1\le d\le f$ に対して
\[
h^{(d)}(\vec x)=h^{(d)}(\vec y),
\qquad
\D_L^{(d)}(\vec x)\cap\D_L^{(d)}(\vec y)\ne\varnothing
\]
ならば
\[
\D_L^{(d)}(\vec x)=\D_L^{(d)}(\vec y)
\]
である。
同値に、有限型が等しく、かつ一つの受理文脈を共有するなら、あらゆるタプル文脈で交換可能である。

\item \textbf{意味論的 clause インターフェース：}
任意の unsafe な意味論的対
$(\vec x,\vec y)\in\mathfrak U_d(L)$, $d\le f$
は少なくとも一つの座標で分離される。
すなわち、ある $i$ について
$h(x_i)\ne h(y_i)$
である。

\item \textbf{合同インターフェース：}
核合同 $\ker h$ は $L$ に対して $f$-分離的である。

\item \textbf{有限統語／関係的インターフェース：}
標準関係的射 $\rho_h:T\relto M$ は点付き統語モノイド $(T,P)$ に対して $f$-分離的である。
\end{enumerate}
したがって、SCL、語／文脈、合同、および標準関係的射による定式化は、一つの安全性述語を witness レベルで表現したものであり、段階的な緩和でも、別々の不変量でもない。
\end{theorem}

\begin{proof}
\textup{(i)} と \textup{(ii)} の同値性は、主 intent の対応関係そのものである。
$\gamma_d(\vec x)$ の intent は $\D_L^{(d)}(\vec x)$ である。
したがって二つの主概念が相異なりかつ隣接することは、それらの分布が重なり、かつ等しくないことと正確に同値である。
そのような隣接概念の生成タプルを選び、その代表元が成分ごとに等しい $h$-型をもつことは、\textup{(ii)} の含意が偽であることと正確に一致する。

\textup{(ii)} と \textup{(iii)} の同値性も、同じ内容を clause 形式で述べただけである。
unsafe pair とは定義により、共通受理文脈をもち、かつ完全分布が異なる対である。
したがって安全性とは、unsafe pair がすべての座標で等しい $h$-値をもたないことにほかならない。
さらに
$x\equiv_{\ker h}y$
であることと $h(x)=h(y)$ であることは同値なので、これは $\ker h$ の $f$-分離条件そのものであり、
\textup{(iii)}$\Longleftrightarrow$\textup{(iv)}
が従う。

残るのは、意味論的インターフェースと有限統語インターフェースの比較である。
\textup{(iii)} が成り立つとし、ある
$(\mathbf s,\mathbf t)\in\mathcal U_d(T,P)$, $d\le f$
が $\rho_h$ によって分離されないと仮定する。
各座標について
\[
m_i\in\rho_h(s_i)\cap\rho_h(t_i)
\]
を選び、さらに
\[
\eta(x_i)=s_i,
\quad
\eta(y_i)=t_i,
\quad
h(x_i)=h(y_i)=m_i
\]
となる語 $x_i,y_i$ を選ぶ。
補題~\ref{alg:lem:profile} により、語タプル $\vec x,\vec y$ は unsafe な意味論的対をなすが、どの座標も $h$ によって分離されない。
これは \textup{(iii)} に反する。
したがって $\rho_h$ は $f$-分離的である。

逆に \textup{(v)} が成り立つとし、
$(\vec x,\vec y)\in\mathfrak U_d(L)$, $d\le f$
を取る。
\[
\mathbf s=\eta^{(d)}(\vec x),
\qquad
\mathbf t=\eta^{(d)}(\vec y)
\]
と置く。
補題~\ref{alg:lem:profile} より
$(\mathbf s,\mathbf t)\in\mathcal U_d(T,P)$
である。
もしすべての $i$ について $h(x_i)=h(y_i)$ なら、この共通値は各座標で
$\rho_h(s_i)\cap\rho_h(t_i)$
に属し、$\rho_h$ の $f$-分離性に反する。
よってある座標が分離され、\textup{(iii)} が従う。
\end{proof}

\begin{table}[t]
\centering
\caption{型付き安全性の対応表。最初の4行は固定された圧縮 witness $h$ に対して同値であり、最後の2行は任意の関係的 witness $\rho$ に対して同値である。}
\label{tab:typed-safety-dictionary}
\small
\begin{tabular}{@{}>{\raggedright\arraybackslash}p{0.27\linewidth}
                  >{\raggedright\arraybackslash}p{0.28\linewidth}
                  >{\raggedright\arraybackslash}p{0.34\linewidth}@{}}
\toprule
\textbf{台集合} & \textbf{witness} & \textbf{実行可能性述語}\\
\midrule
主 SCL／語 & $h:\Sigma^*\to M$ & $\operatorname{Safe}_f(h;L)$\\
語の衝突 & $h:\Sigma^*\to M$ & $\mathfrak U_d(L)$ の座標ごとの分離\\
自由モノイド & $\ker h$ & $L$ に対して $f$-分離的\\
統語モノイド & 標準 $\rho_h:T\relto M$ & $(T,P)$ に対して $f$-分離的\\
\midrule
統語モノイド & 任意の $\rho:T\relto M$ & $(T,P)$ に対して $f$-分離的\\
統語モノイド & 重なり $\tau_\rho$ & $(T,P)$ に対して $f$-safe\\
\bottomrule
\end{tabular}
\end{table}

\begin{remark}[インターフェースの witness レベルでの範囲]
\label{alg:rem:witness-scope}
定理~\ref{alg:thm:safety-interface} の関係的な主張は、固定した圧縮写像 $h$ が誘導する標準関係 $\rho_h$ だけに関するものである。
任意の関係的射が、ある準同型の標準関係そのものになると言っているわけではない。
以下の補題~\ref{alg:lem:selector} は別の自由性 argument であり、任意の関係的射から、その標準 fiber が元の fiber に含まれる関数的 witness を取り出す。
この区別によって、witness レベルのインターフェースが、定理~\ref{alg:thm:relational-characterization} の最小終域特徴づけへ昇格する。
\end{remark}

\subsection{自由モノイド合同による特徴づけ}

上の witness レベル同値性から、自由モノイド上の最適化命題が直ちに得られる。

\begin{theorem}[有限指数合同による特徴づけ]
\label{alg:thm:free-congruence}
任意の言語 $L\subseteq\Sigma^*$ と任意の $f\ge1$ に対して
\[
\cmp_f(L)
=
\min\left\{
[\Sigma^*:\theta]:
\begin{array}{l}
\theta\text{ は有限指数モノイド合同であり、}\\
L\text{ に対して }f\text{-分離的}
\end{array}
\right\}.
\]
以下も含め、そのような有限指数合同が存在しない場合、最小値は $\infty$ と解釈する。
\end{theorem}

\begin{proof}
$h:\Sigma^*\to M$ を有限の witness 圧縮写像とし、$M$ を $\operatorname{im}(h)$ に置き換える。
Safety Interface Theorem により $\ker h$ は $f$-分離的であり、第一同型定理から
\[
[\Sigma^*:\ker h]=|\operatorname{im}(h)|
\]
である。
したがって右辺の最小値は $\cmp_f(L)$ 以下である。

逆に、$\theta$ を有限指数の $f$-分離合同とし、
$q_\theta:\Sigma^*\to\Sigma^*/\theta$
を商射とする。
その核は正確に $\theta$ なので、定理~\ref{alg:thm:safety-interface} の合同インターフェース
\textup{(iv)}$\Rightarrow$\textup{(i)}
により、$q_\theta$ はアリティ $f$ まで safe である。
したがって
\[
\cmp_f(L)
\le |\Sigma^*/\theta|
=[\Sigma^*:\theta].
\]
最小値を取れば等式が従う。
\end{proof}

認識合同は実行可能解の制限されたクラスをなす。
有限指数モノイド合同が $L$ を saturated にするなら、その商射は $L$ を認識し、命題~\ref{conf:prop:recognition-relaxation} により任意の有限アリティで safe である。
したがって、有限指数 $f$-分離合同全体を $\mathsf{Sep}_f(L)$ と書けば
\[
\mathsf{Rec}(L)
\subseteq
\bigcap_{f\ge1}\mathsf{Sep}_f(L),
\qquad
\mathsf{Sep}_{f+1}(L)\subseteq\mathsf{Sep}_f(L).
\]
正則言語 $L$ では、定理~\ref{alg:thm:free-congruence} で合同指数を最小化することにより、入れ子状の緩和鎖
\[
\cmp_1(L)\le\cmp_2(L)\le\cdots\le |T_L|
\]
を得る。
したがってアリティを上げることは、順次より強い分離制約を課すが、一般には厳密な統語認識まで強制するわけではない。

\begin{remark}[clause 形式]
unsafe なアリティ $d$ の対 $(\vec x,\vec y)$ に対し、
$\vec x=(x_1,\ldots,x_d)$, $\vec y=(y_1,\ldots,y_d)$ と書く。
分離要件は clause
\[
\bigvee_{i=1}^d(x_i\not\equiv_\theta y_i)
\]
である。
したがって $\cmp_f(L)$ は、幅 $f$ 以下のすべての unsafe 分離節を満たす有限モノイド合同の最小指数である。
$f=1$ では、各制約は単に必要分離
$x\not\equiv_\theta y$
となる。
\end{remark}

\subsection{関係的射による特徴づけ}

\begin{lemma}[関係的射に対する自由モノイド selector]
\label{alg:lem:selector}
$\rho:T\relto M$ をモノイド関係的射とする。
各文字 $a\in\Sigma$ に対し
$m_a\in\rho(\eta(a))$
を一つ選ぶ。
すると一意なモノイド準同型
\[
h_\rho:\Sigma^*\to M,
\qquad
h_\rho(a)=m_a\quad(a\in\Sigma)
\]
が存在し、任意の $w\in\Sigma^*$ に対して
\[
h_\rho(w)\in\rho(\eta(w))
\]
を満たす。
したがって
\[
\rho_{h_\rho}(t)\subseteq\rho(t)
\qquad(t\in T).
\]
\end{lemma}

\begin{proof}
これは関係的射に対する標準的な自由モノイド選択 argument であるが、以下の像サイズ比較で fiber の包含を明示的に用いるため証明を記す。
$\Sigma^*$ は $\Sigma$ 上の自由モノイドなので、一意な拡張が存在する。

$|w|$ に関する帰納法で
$h_\rho(w)\in\rho(\eta(w))$
を示す。
$w=\varepsilon$ では
\[
h_\rho(\varepsilon)=1_M
\in\rho(1_T)
=\rho(\eta(\varepsilon)).
\]
$w=va$ とし、帰納法の仮定と選んだ文字値から
\[
h_\rho(v)\in\rho(\eta(v)),
\qquad
h_\rho(a)=m_a\in\rho(\eta(a))
\]
を得る。
関係的射の積法則より
\[
h_\rho(w)
=h_\rho(v)h_\rho(a)
\in\rho(\eta(v)\eta(a))
=\rho(\eta(w)).
\]
$m\in\rho_{h_\rho}(t)$ なら、$\eta(w)=t$ を満たすある語 $w$ について $m=h_\rho(w)$ である。
よって $m\in\rho(t)$ であり、
$\rho_{h_\rho}(t)\subseteq\rho(t)$
が従う。
\end{proof}

\begin{theorem}[関係的射による特徴づけ]
\label{alg:thm:relational-characterization}
統語点付きモノイド $(T,P)$ をもつ任意の正則言語 $L$ に対して
\[
\cmp_f(L)
=
\min\left\{
|M|:
\begin{array}{l}
M\text{ は有限モノイドであり、}\\
\rho:T\relto M\text{ は }f\text{-分離的な関係的射}
\end{array}
\right\}.
\]
\end{theorem}

\begin{proof}
右辺の実行可能集合は空でない。
実際、恒等準同型 $T\to T$ を関数的な関係的射とみなせば、これは $f$-分離的である。
unsafe pair $(\mathbf s,\mathbf t)$ について、すべての座標が等しければ profile も同一になってしまうため、少なくとも一座標は異なるからである。

まず、関係的射が未使用の終域値をもつ必要はないことに注意する。
$\rho:T\relto M$ に対し
\[
M':=\pi_2(\operatorname{Gr}(\rho))
=\bigcup_{t\in T}\rho(t)
\]
は $M$ の部分モノイドである。
$1_M$ を含み、関係的射の積法則により積について閉じているからである。
$M$ を $M'$ に置き換えても fiber は何も変わらず、したがって分離性も保存される。
よって右辺の最小値では、$\rho$ が実際に用いる値だけからなる終域に常に制限してよい。

$h:\Sigma^*\to M$ を最小 witness 圧縮写像とし、終域を $\operatorname{im}(h)$ に制限する。
定理~\ref{alg:thm:safety-interface} により、標準関係
$\rho_h:T\relto\operatorname{im}(h)$
は $f$-分離的である。
したがって
\[
\min_\rho |M|\le\cmp_f(L).
\]
逆に、$\rho:T\relto M$ を任意の有限 $f$-分離関係的射とする。
補題~\ref{alg:lem:selector} のように生成元の値を選び、
$h_\rho:\Sigma^*\to M$
を作る。
各 $t$ について
$\rho_{h_\rho}(t)\subseteq\rho(t)$
なので、任意の fiber の非交差
\[
\rho(s_i)\cap\rho(t_i)=\varnothing
\]
は
\[
\rho_{h_\rho}(s_i)\cap\rho_{h_\rho}(t_i)=\varnothing
\]
を含意する。
したがって $\rho_{h_\rho}$ も $f$-分離的である。
定理~\ref{alg:thm:safety-interface} により、$h_\rho$ は tuple substitutability の witness となる。

その終域を $\operatorname{im}(h_\rho)$ に制限すれば
\[
\cmp_f(L)
\le |\operatorname{im}(h_\rho)|
\le |M|.
\]
関係的射全体で最小値を取れば逆向きの不等式が得られる。
\end{proof}

\begin{corollary}[統語不変性]
\label{alg:cor:syntactic-invariance}
$L$ と $K$ を、異なる有限アルファベット上でもよい正則言語とする。
それらの点付き統語モノイドが同型ならば、任意の有限 $f\ge1$ に対して
\[
\cmp_f(L)=\cmp_f(K)
\]
である。
したがって圧縮列全体は、有限点付き統語モノイド $(T_L,P_L)$ の同型に関する不変量である。
\end{corollary}

\begin{proof}
定理~\ref{alg:thm:relational-characterization} の右辺は点付きモノイドだけに依存する。
点付きモノイド同型は unsafe 統語 profile を移送し、さらに合成によって $f$-分離関係的射を両方向へ移送するが、終域サイズは変化させない。
\end{proof}

この特徴づけは有限点付き統語モノイド上で有効的である。
恒等射から $\cmp_f(L)\le|T_L|$ が得られ、アリティ $f$ までの unsafe 集合は有限で計算可能であり、そのサイズ以下の有限 target monoid と関係的射を列挙できる。
ただし効率性を主張しているわけではない。
selector lemma は、多値 fiber が整合性の問題を生じない理由も説明する。
生成元上で値を選べば、自由性により一意に拡張され、乗法性によって関係 fiber の内部に留まり、fiber の交わりを縮小することはあっても増やすことはない。
したがって任意の分離関係的射は、それ以下の像サイズをもつ関数的圧縮 witness を内部に含んでいる。

\subsection{有限側の基本上界：可換な場合の崩壊}

\begin{theorem}[可換統語モノイドにおけるアリティ崩壊]
\label{group:thm:commutative-collapse}
$L$ を正則とし、その統語モノイド $T_L$ が可換であるとする。
任意の有限圧縮写像 $h$ と任意の $f\ge1$ に対して
\[
L\text{ が }(1,h)\text{-substitutable}
\iff
L\text{ が }(f,h)\text{-tuple-substitutable}
\]
である。
したがって
\[
\cmp_1(L)=\cmp_2(L)=\cdots.
\]
\end{theorem}

\begin{proof}
順方向だけ示せばよい。
$\vec x=(x_1,\ldots,x_d)$ と $\vec y=(y_1,\ldots,y_d)$, $d\le f$ が成分ごとに等しい $h$-型をもち、受理タプル文脈
\[
E=u_0\blank_1u_1\cdots\blank_du_d
\]
を共有するとする。
\[
X=x_1\cdots x_d,
\qquad
Y=y_1\cdots y_d,
\qquad
U=u_0\cdots u_d
\]
と置く。
成分ごとの圧縮値の等しさから $h(X)=h(Y)$ である。
統語モノイドが可換なので
\[
\eta(E[\vec x])=\eta(UX),
\qquad
\eta(E[\vec y])=\eta(UY).
\]
したがって $X,Y$ は受理単項文脈 $U\blank_1$ を共有する。
単項 $h$-substitutability により
$\D_L(X)=\D_L(Y)$
である。

別の任意のタプル文脈
\[
F=v_0\blank_1v_1\cdots\blank_dv_d
\]
を取り、$V=v_0\cdots v_d$ と置く。
再び可換性から
\[
F[\vec x]\in L
\iff VX\in L
\iff VY\in L
\iff F[\vec y]\in L.
\]
したがってタプル分布は等しい。
\end{proof}

\section{\texorpdfstring{$\mathbf V_{ab}$}{Vab} 内のアリティ増幅}
\label{sec:vab-arity}

中心指数定理、定理~\ref{group:thm:center-index} は、最初のアリティ上昇だけですでに大きな跳びが起こりうることを示す。
ここではさらに、統語モノイドを一つの固定された古典的擬多様体の内部に保ったまま、任意の連続アリティ境界でその跳びを任意に大きくできることを示す。

\[
M_{ab}:=\operatorname{Synt}(\{ab\})
=\{1,a,b,ab,0\},
\qquad
\mathbf V_{ab}:=\langle M_{ab}\rangle
\]
とする。
モノイド $M_{ab}$ は、最小非可換 variety の Margolis--Pin 分類に現れる標準的な非周期生成子の一つである~\cite{MargolisPin1984}。

語 $w\in\Delta^*$ が\emph{文字相異なる（letter-distinct）}とは、$\Delta$ の各文字が $w$ に高々1回しか現れないことをいう。
言語のすべての語が letter-distinct であるとき、その言語も letter-distinct と呼ぶ。
Margolis と Pin はこのような語を\emph{multilinear} と呼び、彼らの Proposition~1.1 では、variety に $M_{ab}$ が含まれることを、対応する言語 variety にすべての multilinear singleton language が含まれることと同値に特徴づけている~\cite[Proposition~1.1]{MargolisPin1984}。
したがって次の補題は、その結果の有限和に関する帰結としても得られる。
ただし以下で $T_L\in\mathbf V_{ab}$ という形を繰り返し用いるため、明示的な認識 argument を記しておく。

\begin{lemma}[letter-distinct な有限言語は $\mathbf V_{ab}$ に属する]
\label{vab:lem:letterdistinct}
$\Delta$ を有限とする。
$L\subseteq\Delta^*$ が有限で、$L$ のすべての語の文字が相異なるならば
\[
T_L\in\mathbf V_{ab}.
\]
\end{lemma}

\begin{proof}
$w=c_1\cdots c_m$ の文字がすべて相異なるとする。
各 $i$ に対し、$c_i\mapsto ab$ とし、その他のすべての文字を $M_{ab}$ の $1$ へ送る。
像が $ab$ であることを要求すれば、$c_i$ がちょうど1回現れることが強制される。
各 $i<m$ に対しては、$c_i\mapsto a$, $c_{i+1}\mapsto b$ とし、その他の文字をすべて $1$ へ送る。
各文字がちょうど1回現れるという条件のもとで像が $ab$ であることを要求すると、$c_i$ が $c_{i+1}$ より前に現れることが強制される。
さらに別の factor で、$w$ に現れない文字を $0$ へ送る。
したがって有限直積によって singleton $\{w\}$ を正確に認識でき、その統語モノイドは $M_{ab}$ の有限冪の divisor である。

有限和 $L$ に対しては、これら singleton recognizer の直積を取り、少なくとも一つの singleton 座標が受理状態にある積状態を受理集合とする。
各座標状態だけで対応する singleton の受理集合への所属が決まるので、これは積モノイドの well-defined な部分集合である。
したがって積射は正確に $L$ を認識し、$T_L$ は $M_{ab}$ の有限冪の divisor となる。
\end{proof}

$d\ge1$ を固定し、$n=d+1$ と置く。
$1\le j\le R$ に対して互いにすべて異なる文字
\[
a_{j,1},\ldots,a_{j,n},
\qquad
c_{j,1},\ldots,c_{j,n-1}
\]
を導入する。
\[
B_j=a_{j,1}\cdots a_{j,n},
\]
\[
D_j=a_{j,1}c_{j,1}a_{j,2}c_{j,2}\cdots c_{j,n-1}a_{j,n}
\]
とし、
\[
\widehat L_{d,R}:=\{B_j,D_j:1\le j\le R\}
\]
と定める。
受理語はすべて letter-distinct なので
$T_{\widehat L_{d,R}}\in\mathbf V_{ab}$
である。

\begin{lemma}[アリティ $d$ までの圧縮は有界]
\label{vab:lem:low}
$n=d+1$ とする。
$R$ に依存しない圧縮写像 $h_d$ が存在し、$\widehat L_{d,R}$ は $(d,h_d)$-tuple-substitutable であり、
\[
|\operatorname{im}(h_d)|
\le C_d
:=\frac{5d^2+9d+8}{2}
\]
を満たす。
特に
\[
\cmp_d(\widehat L_{d,R})\le C_d.
\]
\end{lemma}

\begin{proof}
\[
\Gamma_d=\{A_1,\ldots,A_n,C_1,\ldots,C_{n-1}\}
\]
とし、
$\rho_d:\Sigma_{d,R}^*\to\Gamma_d^*$
を
\[
a_{j,i}\mapsto A_i,
\qquad
c_{j,i}\mapsto C_i
\]
で定める。
二つの受理 role template は
\[
\beta=A_1\cdots A_n,
\qquad
\delta=A_1C_1A_2C_2\cdots C_{n-1}A_n
\]
である。
空 factor を含め
\[
F_d=\operatorname{Fact}(\{\beta,\delta\})
\]
と置く。
$\Gamma_d^*\setminus F_d$ は両側イデアルなので、Rees factor
\[
Q_d=F_d\cup\{0\}
\]
は有限モノイドである。
すなわち、factor 同士の積は、その連接が再び factor なら連接そのもの、そうでなければ $0$ とする。
商射 $q_d:\Gamma_d^*\to Q_d$ を取り、
$h_d=q_d\circ\rho_d$
と置く。
最長 template の長さは $2n-1$ である。
$\beta$ と $\delta$ の非空 factor の個数を粗く足し合わせれば
\[
|Q_d|
\le
2+\frac{n(n+1)}2+n(2n-1)
=\frac{5d^2+9d+8}{2}
=C_d.
\]
$q_d$ は $F_d$ 上で単射である。

$1\le r\le d=n-1$ とし、
$\mathbf x=(x_1,\ldots,x_r)$ と
$\mathbf y=(y_1,\ldots,y_r)$
が成分ごとに等しい $h_d$-型をもち、受理タプル文脈 $E$ を共有するとする。
受理された充填に現れる各成分は受理語の factor なので、その role word は $F_d$ に属する。
$h_d$-型の等しさと $F_d$ 上での単射性により
\[
\rho_d(x_i)=\rho_d(y_i)
\qquad(1\le i\le r)
\]
である。
二つの受理充填に $\rho_d$ を適用すると、role template は同じであり、$\beta$ または $\delta$ のいずれかである。

二つの充填が同じ copy $j$ に属するなら、$B_j$ または $D_j$ の factor は role word によって一意に決まるので、すべての $i$ について $x_i=y_i$ である。
したがって完全なタプル分布は等しい。

以下、二つの充填が異なる copy $j\ne k$ に属すると仮定する。
すべてのタプル成分が空なら $\mathbf x=\mathbf y$ であり、二つの完成語も一致するので、異なる copy に属することはできない。
したがって少なくとも一つのタプル成分は非空である。
copy ごとのアルファベットは互いに素なので、共有受理文脈に固定 terminal が一つでもあれば、完成語の一方が二つの copy alphabet を混ぜてしまう。
ゆえに $E$ は固定 terminal material をもたない。
したがって非空タプル成分は、hole の出現順序に従って受理語全体を分割し、二つの充填は同じ role template をもつ。

まず template が $\delta$ であるとする。
タプル成分の総長は $|\delta|=2n-1$ であり、これは受理語の最大長である。
任意のタプル文脈 $E'$ を取る。
もし
$E'[\mathbf x]\in\widehat L_{d,R}$
なら、挿入されるタプルだけですでに総長 $2n-1$ なので、$E'$ は固定 terminal material を含めない。
したがって充填語は、固定された hole 順でタプル成分を連接したものにすぎない。
$x_i$ は copy-$j$ のアルファベットだけを使うため、この語は $D_j$ でなければならない。
$\mathbf y$ の対応する等しい role block を連接すれば $D_k$ になるので、
$E'[\mathbf y]\in\widehat L_{d,R}$
である。
$j$ と $k$ を入れ替えれば逆向きも従う。
よって任意の $E'$ に対して
\[
E'[\mathbf x]\in\widehat L_{d,R}
\quad\Longleftrightarrow\quad
E'[\mathbf y]\in\widehat L_{d,R}
\]
であり、二つのタプル分布は等しい。

最後に共通 template が $\beta$ であるとする。
非空タプル成分は $A_1\cdots A_n$ を高々 $r\le n-1$ 個の連続 role block に分割する。
したがって少なくとも一つの block は隣接対 $A_iA_{i+1}$ を含む。
再び任意の $E'$ を取り、まず
$E'[\mathbf x]\in\widehat L_{d,R}$
とする。
挿入文字はすべて copy $j$ に属するので、受理充填が別 copy の固定文字を混ぜることはできない。
target が $B_j$ なら、その長さは挿入総長 $n$ とすでに等しいので、$E'$ は固定 terminal material をもたない。
そこで copy $j$ を copy $k$ に置き換えれば $B_k$ となり、$E'[\mathbf y]$ も受理される。

copy-$j$ alphabet 上のもう一つの可能な target は $D_j$ だが、その role template は $\delta$ である。
しかし $D_j$ では marker $C_i$ が $A_i$ と $A_{i+1}$ の間に厳密に置かれる一方、この二つの role は一つのタプル成分の内部に入っている。
固定文脈 material を成分内部に挿入することはできないので、$D_j$ は形成できない。
したがって $\mathbf x$ を受理する任意の文脈は $\mathbf y$ も受理する。
$j$ と $k$ を入れ替えれば逆向きも従う。
よって再び任意の $E'$ に対して
\[
E'[\mathbf x]\in\widehat L_{d,R}
\quad\Longleftrightarrow\quad
E'[\mathbf y]\in\widehat L_{d,R}.
\]

すべての場合に完全なタプル分布が一致するので、$h_d$ はアリティ $d$ までの witness である。
\end{proof}

\begin{lemma}[次アリティは非有界]
\label{vab:lem:high}
$n=d+1$ に対して
\[
\cmp_n(\widehat L_{d,R})
\ge
\left\lceil R^{1/n}\right\rceil.
\]
\end{lemma}

\begin{proof}
\[
\mathbf a_j=(a_{j,1},\ldots,a_{j,n})
\]
と置く。
隣接 hole 文脈
$\blank_1\cdots\blank_n$
はすべての $\mathbf a_j$ を受理し、$B_j$ を生成する。
一方、文脈
\[
\blank_1c_{j,1}\blank_2c_{j,2}\cdots c_{j,n-1}\blank_n
\]
は $\mathbf a_j$ を受理して $D_j$ を生成するが、$k\ne j$ のすべての $\mathbf a_k$ を拒否する。
したがって $R$ 個のタプルは pairwise unsafe である。
像に $m$ 個の値をもつ witness を取ると、成分ごとの型は高々 $m^n$ 個しかなく、それらはすべて異ならなければならない。
よって $R\le m^n$ である。
\end{proof}

\begin{theorem}[$\mathbf V_{ab}$ 内の全境界増幅]
\label{vab:thm:amplification}
任意の固定 $d\ge1$ に対して
\[
\sup_{\,T_L\in\mathbf V_{ab}}
\frac{\cmp_{d+1}(L)}{\cmp_d(L)}
=\infty.
\]
witness は、語長が高々 $2d+1$ の有限 letter-distinct 言語として取れる。
さらに、ある整数 $m_d\ge1$ と、そのような言語の列 $L_j$ で
$T_{L_j}\in\mathbf V_{ab}$ を満たし、
\[
\cmp_d(L_j)=m_d\quad\text{for all }j,
\qquad
\cmp_{d+1}(L_j)\longrightarrow\infty
\]
となるものが存在する。
\end{theorem}

\begin{proof}
$\widehat L_{d,R}$ に対し、補題~\ref{vab:lem:low} と~\ref{vab:lem:high} から
\[
\frac{\cmp_{d+1}(\widehat L_{d,R})}
     {\cmp_d(\widehat L_{d,R})}
\ge
\frac{\lceil R^{1/(d+1)}\rceil}{C_d}
\]
を得る。
これは $R\to\infty$ で無限大へ発散する。
同じ補題から
\[
1\le\cmp_d(\widehat L_{d,R})\le C_d,
\qquad
\cmp_{d+1}(\widehat L_{d,R})
\ge\lceil R^{1/(d+1)}\rceil
\]
である。
したがって $\cmp_d(\widehat L_{d,R})$ は有限個の整数値しか取らない。
そのうち一つは無限個の $R$ に対して現れるので、その部分列に沿って低アリティ値は一定のまま、次アリティ値だけが発散する。
\end{proof}

\subsection{\texorpdfstring{$\mathbf V_{ab}$}{Vab} 内の有限初期部分普遍性}

\begin{lemma}[互いに素な和に対する上下界]
\label{vab:lem:union}
$L_i\subseteq\Sigma_i^*$ を、互いに素な有限アルファベット上の言語とする。
すべての $i$ について $\varepsilon\notin L_i$ と仮定し、
\[
L=\bigcup_{i=1}^sL_i
\]
と置く。
このとき
\[
\max_i\cmp_f(L_i)\le\cmp_f(L).
\]
また、$h_i:\Sigma_i^*\to M_i$ が $L_i$ の witness ならば、$L$ は高々
\[
2^s\prod_{i=1}^s|M_i|
\]
個の像元をもつ witness をもつ。
\end{lemma}

\begin{proof}
下界について、$h$ を $L$ の witness とし $i$ を固定する。
アルファベットが互いに素で、どの $L_j$ も空語を含まないので
\[
L_i=L\cap\Sigma_i^*.
\]
$L_i$ の任意の unsafe tuple pair は、$\Sigma_i$ 上の文脈によって受理され区別されるので、$L$ においても unsafe のままである。
したがって $h$ の $\Sigma_i^*$ への制限は $L_i$ の witness であり、
\[
\cmp_f(L_i)\le|\operatorname{im}(h)|
\]
を得る。

上界について、$\pi_i:\Sigma^*\to\Sigma_i^*$ を $\Sigma_i$ 外の文字をすべて消去する写像とし、
\[
\operatorname{supp}:\Sigma^*\longrightarrow(2^{[s]},\cup)
\]
を、語に現れる component alphabet の集合を記録する写像とする。
\[
H=(\operatorname{supp},h_1\circ\pi_1,\ldots,h_s\circ\pi_s)
\]
と定める。
その像のサイズは高々
$2^s\prod_i|M_i|$
である。

成分ごとに $H$-等しいタプル $\mathbf x,\mathbf y$ が受理タプル文脈 $E$ を共有するとする。
すべてのタプル成分が空なら $\mathbf x=\mathbf y$ なので、少なくとも一成分が非空とする。
$E[\mathbf x]\in L$ なので、アルファベットが互いに素であることから、それは一意な $L_j$ に属する。
したがって $E$ のすべての固定 terminal と、すべての非空 $x_k$ は $\Sigma_j$ の文字だけを使う。
support 座標の等しさから、対応する各 $y_k$ も正確に同じ alphabet support をもち、したがって $E[\mathbf y]$ も $L_j$ に属する。
さらに成分ごとの $H$-等値性により
\[
h_j(\pi_j(x_k))=h_j(\pi_j(y_k))
\qquad\text{for every coordinate }k,
\]
しかもここでは $\pi_j(x_k)=x_k$, $\pi_j(y_k)=y_k$ である。
$h_j$ は $L_j$ の witness なので、二つのタプルは $\Sigma_j$ 上のすべての名前付き文脈に関して同じ分布をもつ。

残るのは全 union alphabet 上の文脈の比較である。
二つのタプルのどちらか一方を受理する任意の文脈 $F$ は、再びそれをある $L_k$ へ完成させなければならない。
非空タプル成分の存在と support の等しさにより $k=j$ が強制される。
したがってそのような受理文脈 $F$ の固定 terminal はすべて $\Sigma_j$ に属する。
よって $F$ は実質的に $\Sigma_j$-文脈であり、前段落により二つのタプルに対する受理性は同じである。
複数 component alphabet を混ぜる文脈は、どちらのタプルも受理しない。
ゆえに $L$ における完全タプル分布は等しく、$H$ は $L$ の witness である。
\end{proof}

\begin{lemma}[三角形的貼り合わせ原理]
\label{vab:lem:triangular-gluing}
$m\ge2$ を固定する。
各 $1\le g<m$ に対し、有限アルファベット上の $\varepsilon$-free な正則言語族
$\{K_{g,r}:r\ge1\}$
があるとする。
また整数 $C_g\ge1$ と関数
$B_g:\mathbb N\to\mathbb R_{>0}$
で $B_g(r)\to\infty$ を満たし、任意の $r$ に対して
\[
\cmp_g(K_{g,r})\le C_g,
\qquad
\cmp_{g+1}(K_{g,r})\ge B_g(r)
\tag{TG}
\]
とする。
このとき任意の
$\lambda_1,\ldots,\lambda_{m-1}>0$
に対して、パラメータ $r_1,\ldots,r_{m-1}$ を選び、component alphabet を互いに素になるよう rename し、
\[
K:=\bigcup_{g=1}^{m-1}K_{g,r_g}
\]
とすれば、
\[
\cmp_{f+1}(K)>\lambda_f\cmp_f(K)
\quad\text{かつ}\quad
\cmp_{f+1}(K)>\cmp_f(K)
\qquad(1\le f<m)
\]
を同時に満たすようにできる。
特に、\textup{(TG)} を満たす任意の境界 amplifier の集合を一つの言語へ貼り合わせ、その連続圧縮比を、あらかじめ指定した任意の有限乗法ギャップ閾値より同時に大きくできる。
\end{lemma}

\begin{proof}
固定した component $K_{g,r}$ に対し、厳密な統語認識は全アリティ witness
\[
D_g(r):=|T_{K_{g,r}}|
\]
を与える。
アルファベットの rename は圧縮数を変えないので、選んだ component はすべて pairwise alphabet-disjoint と仮定してよい。

$r_1,\ldots,r_{m-1}$ を再帰的に選ぶ。
$r_1,\ldots,r_{f-1}$ まで選んだ段階で
\[
U_f
:=2^{m-1}
\left(\prod_{g<f}D_g(r_g)\right)
\left(\prod_{g\ge f}C_g\right)
\]
と置く。
重要なのは、$U_f$ がすでに選ばれたパラメータだけに依存することである。
$g\ge f$ の component は、将来のパラメータではなく、一様な低アリティ上界 $C_g$ だけを通じて入る。
$B_f(r)\to\infty$ なので
\[
B_f(r_f)>\max\{1,\lambda_f\}U_f
\]
となる $r_f$ を選べる。
これが再帰を循環させない三角形的特徴である。

$K=\bigcup_{g=1}^{m-1}K_{g,r_g}$ とする。
アリティ $f$ において、すでに選択済みの各 component $g<f$ にはサイズ $D_g(r_g)$ の統語的全アリティ witness を用いる。
$g\ge f$ については条件 \textup{(TG)} がアリティ $g$ まで、したがってアリティ $f$ まで、高々 $C_g$ 値の witness を与える。
補題~\ref{vab:lem:union} の上界側から
\[
\cmp_f(K)\le U_f
\]
を得る。
一方、同補題の下界側と $f$ 番目の amplifier から
\[
\begin{aligned}
\cmp_{f+1}(K)
&\ge\cmp_{f+1}(K_{f,r_f})
\ge B_f(r_f)\\
&>\max\{1,\lambda_f\}U_f
\ge\max\{1,\lambda_f\}\cmp_f(K).
\end{aligned}
\]
これで任意の $1\le f<m$ に対して二つの不等式が同時に示された。
\end{proof}

\begin{theorem}[$\mathbf V_{ab}$ 内の有限初期部分普遍性]
\label{vab:thm:finite-prefix}
$m\ge2$ および $\lambda_1,\ldots,\lambda_{m-1}>0$ とする。
語長がすべて高々 $2m-1$ である有限 letter-distinct 言語 $L$ で
$T_L\in\mathbf V_{ab}$
を満たし、
\[
\cmp_{f+1}(L)>\lambda_f\cmp_f(L)
\qquad(1\le f<m)
\]
となるものが存在する。
特に
\[
\cmp_1(L)<\cmp_2(L)<\cdots<\cmp_m(L).
\]
\end{theorem}

\begin{proof}
各 $1\le g<m$ に対し、amplifier family
\[
K_{g,R}:=\widehat L_{g,R}
\]
へ補題~\ref{vab:lem:low} と~\ref{vab:lem:high} を適用する。
これらの言語は有限かつ $\varepsilon$-free であり、補題~\ref{vab:lem:triangular-gluing} の仮定を
\[
C_g=\frac{5g^2+9g+8}{2},
\qquad
B_g(R)=\left\lceil R^{1/(g+1)}\right\rceil
\]
で満たす。
したがって三角形的貼り合わせ原理により $R_1,\ldots,R_{m-1}$ を選べ、アルファベットを互いに素になるよう rename した後
\[
L=\bigcup_{g=1}^{m-1}\widehat L_{g,R_g}
\]
がすべての要求不等式を満たすようにできる。

各受理語は一つの component に属し letter-distinct である。
その長さは高々
\[
\max_{g<m}(2g+1)=2m-1.
\]
したがって union も有限 letter-distinct 言語であり、補題~\ref{vab:lem:letterdistinct} により
$T_L\in\mathbf V_{ab}$
である。
\end{proof}

\subsection{SCL 圧縮高さ}

\begin{definition}[SCL 圧縮高さ]
\label{ah:def:height}
$\mathbf V$ を有限モノイドの擬多様体とする。
$\Obj_1(L)\cong T_L$ なので
\[
\operatorname{ch}(\mathbf V)
:=
\min\left\{
r\in\mathbb N_{\ge1}:
\begin{array}{l}
T_L\in\mathbf V\text{ を満たすすべての正則言語 }L\text{ について、}\\
\cmp_f(L)=\cmp_r(L)
\text{ がすべての }f\in\mathbb N_{\ge r}\text{ で成り立つ}
\end{array}
\right\}
\]
と定める。
集合が空なら最小値を $\infty$ とする。
ここでも以後この定義でも、正則言語に関する量化は任意の有限アルファベット上の言語全体を走る。
クラス全体で共通のアルファベットを固定するわけではない。
\end{definition}

各固定正則言語 $L$ に対して、列
\[
\cmp_1(L)\le\cmp_2(L)\le\cdots\le|T_L|
\]
は単調非減少かつ有界なので、ある有限アリティで必ず安定する。
したがって $\operatorname{ch}(\mathbf V)=\infty$ とは、統語モノイドが $\mathbf V$ に属するすべての言語に一様に使える単一の安定化アリティが存在しないことを意味する。
一つの言語について圧縮列が発散することを意味するわけではない。

\begin{lemma}[単調性]
\label{ah:lem:monotone}
$\mathbf V\subseteq\mathbf W$ ならば
\[
\operatorname{ch}(\mathbf V)
\le
\operatorname{ch}(\mathbf W).
\]
\end{lemma}

\begin{proof}
$\mathbf V$ に対して量化される言語クラスは、$\mathbf W$ に対して量化される言語クラスに含まれる。
\end{proof}

\begin{corollary}[二つの基本高さ境界]
\label{height:cor:core-bounds}
有限可換モノイドの擬多様体 $\mathbf{Com}$ は
\[
\operatorname{ch}(\mathbf{Com})=1
\]
を満たす。
一方、
\[
\operatorname{ch}(\mathbf V_{ab})=\infty.
\]
\end{corollary}

\begin{proof}
第一の主張は定理~\ref{group:thm:commutative-collapse} である。
第二は定理~\ref{vab:thm:amplification} から従う。
$\mathbf V_{ab}$ 内ではすべてのアリティ境界が非有界なギャップをもつため、この擬多様体全体に一様に働く有限安定化アリティは存在しない。
\end{proof}

% ===== 日本語化済み：冒頭〜第5節末。以下は未翻訳の原文英語（続き作業用） =====
\section{完全正則モノイドはアリティ2で飽和する}
\label{sec:cr-saturation}

\paragraph{証明戦略.}
補題~\ref{alg:lem:relational-overlap} により、関係的射に対するアリティ問題は、対応する compatible tolerance に対する問題へ帰着される。この tolerance は推移的である必要はないが、2-safe 性は、各受理 completely simple 成分の内部に強い剛性を強制する。関係づけられた元は、その成分内で一意な中心平行移動だけ異なり、これらの局所的な差異は準同型
\(z_\beta:\Gamma_\beta\to N_\beta\le Z(G_\beta)\)
へと組み上がる。共通の受理文脈は座標ごとの欠陥の積を自明にし、文脈的因子分解はその恒等式を他のすべてのタプル文脈でも保存する。これにより、2-safe 性がすべての有限アリティへ引き上げられる。

$(S,P)$ を有限の統語的点付きモノイドとする。以下、統語分離性
\[
 r\ne s
 \quad\Longrightarrow\quad
 \text{$arb,asb$ のちょうど一方だけが $P$ に属するような $a,b\in S$ が存在する。}
\tag{Syn}
\]
を繰り返し用いる。これは、全射な統語射を通して移された、異なる統語類の定義的な分離性である。

Table~\ref{tab:typed-safety-dictionary} の型付き安全性対応表を用いる。特に、補題~\ref{alg:lem:relational-overlap} は任意の関係的 witness \(\rho\) をその重なり tolerance \(\tau_\rho\) へ変換し、\(\rho\) の $f$-分離性と \(\tau_\rho\) の $f$-safe 性が正確に同値であることを与える。したがって、完全正則な統語的点付きモノイド上の任意の $2$-safe compatible tolerance が、すべての有限アリティで safe であることを示せば十分であり、実際これはわずかに強い主張である。推移性は仮定せず、証明中でも用いない。古典的な用語では、このような関係は tolerance と呼ばれる
\cite{ZelinkaTolerance1975,Pondelicek1978,KumaresanTolerance1984}。また直積二乗の言葉では、これは $S\times S$ の対称な対角部分半群である~\cite{BarberRuskucDirectSquare,BarberRuskucDiagonal2026}。

以下では補題~\ref{alg:lem:profile} の有限統語 profile を直接用いる。したがって「タプル文脈」は、具体的な語ではなく $S$ の固定因子をもってよい。統語射の全射性により、そのような各 profile には具体的な語による代表元が存在する。

この節を通して $S$ は完全正則であると仮定する。完全正則半群を completely simple 成分の半束へ分解する Clifford の分解~\cite[Theorem~4.1.3]{Howie1995}、各成分の Rees--Suschkewitsch 表現~\cite[Theorem~3.3.1]{Howie1995}、および正規化された Rees 座標~\cite[Theorem~3.4.2]{Howie1995} を用いる。以下の証明が必要とするのは各受理成分内の正規化 Rees 積だけであり、translational hull の標準形は不要である。

したがって
\(S=\bigsqcup_{\alpha\in Y}S_\alpha,\)
ここで $Y=S/\mathcal D$ は半束、各 $S_\alpha$ は completely simple であり、
\[
S_\alpha S_\beta\subseteq S_{\alpha\beta}.
\tag{CR1}
\]
自然順序
\(\beta\le\alpha \Longleftrightarrow \alpha\beta=\beta\)
を用いる。$x\in S_\alpha$ に対して $\operatorname{supp}(x)=\alpha$ と書く。$S$ はモノイドなので、$1_Y:=\operatorname{supp}(1_S)$ は $Y$ の最大元である。

\begin{lemma}[下位成分への作用とサンドイッチ・イデアル]
\label{cr:lem:sandwich-ideal}
\(c\in S_\alpha\) かつ \(\beta\le\alpha\) とする。このとき
\(
cS_\beta\subseteq S_\beta,
\qquad
S_\beta c\subseteq S_\beta.
\)
したがって \(S_\beta c S_\beta\) は completely simple 半群 \(S_\beta\) の空でない両側イデアルであり、ゆえに
\(S_\beta c S_\beta=S_\beta\)
である。
\end{lemma}

\begin{proof}
\textup{(CR1)} と \(\beta\le\alpha\) より、
\(
S_\alpha S_\beta\subseteq S_{\alpha\beta}=S_\beta,
\qquad
S_\beta S_\alpha\subseteq S_{\beta\alpha}=S_\beta.
\)
したがって \(S_\beta cS_\beta\) は空でない。$z=xcy$ がこれに属し、$r\in S_\beta$ とすると、$rz=(rx)cy$ および $zr=xc(yr)$ もこれに属する。よってこれは $S_\beta$ の両側イデアルであり、completely simple 半群の単純性から結論が従う。
\end{proof}

証明には二つのインターフェースがある。意味論的インターフェースは次の二つの持ち上げ原理からなり、構造的インターフェースは後で証明する中心欠陥準同型である。

\begin{lemma}[safe な重なりからの統語値持ち上げ]
\label{cr:lem:syntactic-lifting}
$1\le d\le f$ とし、$\tau$ を $f$-safe compatible tolerance とする。$1\le i\le d$ に対して $s_i\tau t_i$ であり、二つの $d$-タプルが共通の受理タプル文脈をもつとする。このとき任意の $d$-項タプル文脈 $F$ に対して、
\(F[\mathbf s]=F[\mathbf t]\quad\text{in }S.\)
\end{lemma}

\begin{proof}
二つのタプルは座標ごとに $\tau$-関係にあり、かつ共通の受理 profile をもつので、$f$-safe 性により、それらが $\mathcal U_d(S,P)$ の元をなすことは排除される。したがって有限 profile は等しく、任意の $F$ に対して $F[\mathbf s]\in P$ iff $F[\mathbf t]\in P$ である。もしある $F$ について二つの統語値が異なるなら、\textup{(Syn)} により $a,b\in S$ が存在して、$aF[\mathbf s]b$ と $aF[\mathbf t]b$ のちょうど一方だけが $P$ に属する。すると包み込んだ文脈 $aFb$ が両者の有限 profile を区別してしまい、矛盾する。
\end{proof}

\begin{lemma}[等値化文脈原理]
\label{cr:lem:equalizing-context}
$1\le d\le f$ とし、$\tau$ を $f$-safe compatible tolerance とする。$1\le i\le d$ に対して $s_i\tau t_i$ とする。ある $d$-項タプル文脈 $E_0$ が
\[
E_0[\mathbf s]=E_0[\mathbf t]=c
\]
を満たし、さらに
\(\beta\le\operatorname{supp}(c)\)
を満たす受理成分 $S_\beta$ が存在するとする。このとき
\[
F[\mathbf s]=F[\mathbf t]
\]
が任意の $d$-項タプル文脈 $F$ に対して成り立つ。
\end{lemma}

\begin{proof}
$p\in P\cap S_\beta$ を選ぶ。補題~\ref{cr:lem:sandwich-ideal} により $S_\beta cS_\beta=S_\beta$ だから、$acb=p$ を満たす $a,b\in S_\beta$ が存在する。したがって $aE_0b$ は $\mathbf s$ と $\mathbf t$ の共通受理文脈である。よって補題~\ref{cr:lem:syntactic-lifting} により、任意の $F$ に対して主張の等式が得られる。
\end{proof}

\begin{remark}[アリティ2定理の証明依存関係]
\label{cr:rem:dependencies}
証明の構造的な骨格は短い：
\[
\begin{gathered}
\text{統語持ち上げ + 局所剛性}
\Longrightarrow
\text{自由な中心平行移動}
\Longrightarrow
z_\beta:\Gamma_\beta\to N_\beta
\\[-1pt]
\Longrightarrow
\text{積圧縮}
\Longrightarrow
\text{全アリティ安全性}.
\end{gathered}
\]
より正確には、補題~\ref{cr:lem:syntactic-lifting} と
\ref{cr:lem:equalizing-context} が意味論的インターフェースをなす。それらに補題~\ref{cr:lem:local-kernel}, \ref{cr:lem:local-central}, \ref{cr:lem:index-rigidity} を組み合わせることで局所段階が得られる。その一様な形が、補題~\ref{cr:lem:central-translations} の中心平行移動作用と、命題~\ref{cr:lem:uniform-defect} の準同型である。続いて文脈的因子分解と補題~\ref{cr:lem:product-compression} から定理~\ref{cr:thm:tolerance-saturation} が従う。support の管理が必要になるのは積圧縮の段階だけである。tolerance や対角部分半群の分類、translational-hull の標準形、あるいはアリティ2を超える安全性仮定は一切用いない。
\end{remark}

\subsection{受理 completely simple 成分内の局所剛性}

$P\cap S_\beta\ne\varnothing$ を満たす各 \(\beta\in Y\) について、正規化 Rees 表現
\(S_\beta=\mathcal M[G_\beta;I_\beta,\Lambda_\beta;Q^\beta]\)
を一つ固定する。以下の局所的議論では、\(\beta\) は任意の受理成分を表し、すべての Rees 座標計算はその固定した正規化表現内で行う。
この節では一貫して積の規約
\[
(i,g,\lambda)(j,h,\mu)
=
(i,\,gq^\beta_{\lambda j}h,\,\mu)
\tag{CR2}
\]
を用いる。特別な行と列の添字を $1$ とし、正規化とは
\(
q^\beta_{1i}=q^\beta_{\lambda 1}=1_{G_\beta}
\qquad(i\in I_\beta,\ \lambda\in\Lambda_\beta)
\)
を意味する。以下の Rees 座標計算はすべて \textup{(CR2)} を用いる。
制限
\(\tau_\beta:=\tau\cap(S_\beta\times S_\beta)\)
は $S_\beta$ 上の compatible tolerance、同値に $S_\beta\times S_\beta$ の対称な対角部分半群である。すなわち構成上反射的かつ対称的であり、両立性は上で確認済みである。$\tau$ 自身と同様に、$\tau_\beta$ も推移的である必要はない。

\(H_{11}=\{(1,g,1):g\in G_\beta\}, \qquad e=(1,1,1), \qquad \bar g=(1,g,1)\)
と書く。
\(x=(i,g,\lambda)\in S_\beta\) に対して、\textup{(CR2)} の正規化から
\(exe=\bar g\), \(xe=(i,g,1)\), \(ex=(1,g,\lambda)\)
が得られ、以下でこれらを用いる。
Rees 表現が正規化されているため、$H_{11}$ は $e$ を単位元とし $G_\beta$ と同型な部分群である。

\begin{lemma}[局所核]
\label{cr:lem:local-kernel}
\(
N_\beta:=\{n\in G_\beta:e\tau\bar n\}
\)
と定める。このとき $N_\beta\trianglelefteq G_\beta$ である。
\end{lemma}

\begin{proof}
\(R:=\{(x,y)\in H_{11}\times H_{11}:x\tau y\}\)
とおく。両立性により $R$ は有限群 $H_{11}\times H_{11}$ の部分半群であり、反射性により対角集合、したがって $(e,e)$ を含む。単位元を含む有限群の任意の部分半群は部分群である。実際、$r\in R$ に対し有限性からある $m>0$ で $r^m=(e,e)$ となり、$r^{-1}=r^{m-1}\in R$ である。よって $R$ は部分群である。第一座標が $e$ である fiber はちょうど $\{(e,\bar n):n\in N_\beta\}$ なので、$N_\beta$ は単位元を含み、積と逆元で閉じる。したがって $N_\beta\le G_\beta$ である。

正規性について、$n\in N_\beta$ と $g\in G_\beta$ をとる。両立性により、関係づけられた三つの対
\((\bar g,\bar g),\qquad (e,\bar n),\qquad (\overline{g^{-1}},\overline{g^{-1}})\)
を座標ごとに掛け合わせられる。その積は
\((e,\overline{gng^{-1}})\)
だから、$gng^{-1}\in N_\beta$ である。
\end{proof}

\begin{lemma}[局所中心性]
\label{cr:lem:local-central}
$\tau$ が $2$-safe なら、
\(N_\beta\le Z(G_\beta)\)
である。
\end{lemma}

\begin{proof}
$n\in N_\beta$ とする。補題~\ref{cr:lem:local-kernel} により $n^{-1}\in N_\beta$ でもあるから、タプル
\((e,e)\) と \((\bar n,\overline{n^{-1}})\)
は座標ごとに $\tau$-関係にある。固定因子をもたない二項文脈
$E_0=\blank_1\blank_2$
は両タプル上で同じ値 $e$ を与える。$S_\beta$ は受理成分なので、補題~\ref{cr:lem:equalizing-context} を適用できる。任意の $g\in G_\beta$ に対して
\(F=\blank_1\bar g\blank_2\overline{g^{-1}}\)
で評価すると、
\[
e
=
\bar n\bar g\overline{n^{-1}}\overline{g^{-1}}
=
\overline{ngn^{-1}g^{-1}},
\]
したがって $ngn^{-1}g^{-1}=1$ である。ゆえに $n\in Z(G_\beta)$。
\end{proof}

\begin{lemma}[局所的な行・列剛性]
\label{cr:lem:index-rigidity}
$\tau$ が $2$-safe であり、
\((i,g,\lambda)\tau_\beta(j,h,\mu)\)
とする。このとき $i=j$ かつ $\lambda=\mu$ であり、さらに
\(gh^{-1}\in N_\beta\)
である。
\end{lemma}

\begin{proof}
$x=(i,g,\lambda)$, $y=(j,h,\mu)$ とおく。特別な冪等元 $e$ との両立性により
\begin{equation}
\bar g=exe\;\tau\;eye=\bar h.
\label{eq:cr-compressed-related}
\end{equation}
対称性と両立性から
\(e\tau\overline{h^{-1}g}\)
である。$d_R:=h^{-1}g\in N_\beta$ とおく。するとタプル $(x,e)$ と $(y,\overline{d_R})$ は座標ごとに $\tau$-関係にあり、文脈
\(E_R=e\blank_1\blank_2\)
は共通値 $\bar g$ を与える。補題~\ref{cr:lem:equalizing-context} を適用し、続いて固定因子なしの文脈 $\blank_1\blank_2$ で評価すると
\[
xe=y\overline{d_R}
\]
を得る。正規化 Rees 座標ではこれは
$(i,g,1)=(j,g,1)$
であり、したがって $i=j$。

双対的に、\eqref{eq:cr-compressed-related} における両立性から
$d_L:=gh^{-1}\in N_\beta$ に対して $e\tau\overline{d_L}$ を得る。タプル $(e,x)$ と $(\overline{d_L},y)$ は座標ごとに $\tau$-関係にあり、
\(E_L=\blank_1\blank_2e\)
は共通値 $\bar g$ を与える。補題~\ref{cr:lem:equalizing-context} を適用してから $\blank_1\blank_2$ で評価すると
\[
ex=\overline{d_L}y.
\]
よって $(1,g,\lambda)=(1,g,\mu)$ であり、$\lambda=\mu$。同じ計算によりすでに $gh^{-1}=d_L\in N_\beta$ も得られている。
\end{proof}

したがって各受理 completely simple 成分では、強い局所形
\[
x\tau_\beta y
\quad\Longrightarrow\quad
x,y\text{ は同じ }\mathcal H\text{-類に属し、}
N_\beta\le Z(G_\beta)\text{ の元だけ異なる}
\tag{LR}
\]
が成り立つ。

\subsection{一様な中心欠陥}

前小節で固定した受理成分 $S_\beta$ を考える。
\(\Omega_\beta:=\{a\in S:\beta\le\operatorname{supp}(a)\}\)
とおく。support の積は半束積なので、\(\Omega_\beta\) は $S$ の部分モノイドである。補題~\ref{cr:lem:sandwich-ideal} により、任意の $a\in\Omega_\beta$ は左右どちらから掛けても $S_\beta$ を保存する。

\begin{lemma}[中心平行移動]
\label{cr:lem:central-translations}
\(n\in N_\beta\) に対して
\[
\nu_\beta(n)(i,g,\lambda):=(i,ng,\lambda)
\qquad((i,g,\lambda)\in S_\beta)
\]
と定める。このとき
\(\nu_\beta(mn)=\nu_\beta(m)\circ\nu_\beta(n)\) および
\(\nu_\beta(1)=\operatorname{id}_{S_\beta}\)
が成り立つ。さらに $x,y\in S_\beta$ に対して
\[
\nu_\beta(n)(x)y=x\nu_\beta(n)(y)=\nu_\beta(n)(xy),
\tag{CT}
\]
この作用は自由であり、任意の \(a\in\Omega_\beta\) に対して
\[
a\nu_\beta(n)(x)=\nu_\beta(n)(ax),
\qquad
\nu_\beta(n)(x)a=\nu_\beta(n)(xa).
\tag{CE}
\]
特に、ある $x\in S_\beta$ について
\(\nu_\beta(m)(x)=\nu_\beta(n)(x)\)
ならば $m=n$ である。
\end{lemma}

\begin{proof}
合成則は定義から直ちに従う。
\(x=(i,g,\lambda)\), \(y=(j,h,\mu)\)
と書く。補題~\ref{cr:lem:local-central} により
\(N_\beta\le Z(G_\beta)\) なので、\textup{(CR2)} から
\[
\nu_\beta(n)(x)y
=(i,ngq^\beta_{\lambda j}h,\mu)
=x\nu_\beta(n)(y)
=\nu_\beta(n)(xy),
\]
となり、\textup{(CT)} が得られる。もし
\(\nu_\beta(n)(i,g,\lambda)=(i,g,\lambda)\)
なら $ng=g$ だから $n=1$。したがって作用の自由性と最後の一意性が従う。

\textup{(CE)} を示す。completely simple 半群 $S_\beta$ は $S_\beta=S_\beta^2$ を満たすので、$x=x_1x_2$ と $x_1,x_2\in S_\beta$ に分解できる。$ax_1\in S_\beta$ だから、\textup{(CT)} より
\[
a\nu_\beta(n)(x)
=a\bigl(x_1\nu_\beta(n)(x_2)\bigr)
=(ax_1)\nu_\beta(n)(x_2)
=\nu_\beta(n)(ax).
\]
右側の恒等式は双対的に従う。
\end{proof}

\begin{corollary}[局所中心平行移動形]
\label{cr:cor:local-translation}
$\tau$ が $2$-safe であるとする。$x\tau_\beta y$ ならば、
\(y=\nu_\beta(n)(x)\)
を満たす一意な \(n\in N_\beta\) が存在する。
\end{corollary}

\begin{proof}
\(x=(i,g,\lambda)\), \(y=(j,h,\mu)\)
と書く。補題~\ref{cr:lem:index-rigidity} により $i=j$, $\lambda=\mu$, $gh^{-1}\in N_\beta$ である。したがって $n:=hg^{-1}\in N_\beta$ とおけば $y=\nu_\beta(n)(x)$。一意性は補題~\ref{cr:lem:central-translations} から従う。
\end{proof}

\begin{proposition}[中心欠陥準同型]
\label{cr:lem:uniform-defect}
$\tau$ が $2$-safe であるとし、
\(\Gamma_\beta:=\{(a,b)\in\Omega_\beta\times\Omega_\beta:a\tau b\}\)
とおく。このとき \(\Gamma_\beta\) は $S\times S$ の部分モノイドであり、任意の
\((a,b)\in\Gamma_\beta\) および $x\in S_\beta$ に対して
\[
bx=\nu_\beta\bigl(z_\beta(a,b)\bigr)(ax),
\qquad
xb=\nu_\beta\bigl(z_\beta(a,b)\bigr)(xa)
\tag{UD}
\]
を満たす一意な写像
\(z_\beta:\Gamma_\beta\to N_\beta\)
が存在する。さらに $z_\beta$ はモノイド準同型であり、すなわちすべての
$(a,b),(a',b')\in\Gamma_\beta$ に対して
\begin{equation}
z_\beta(aa',bb')=z_\beta(a,b)z_\beta(a',b')
\label{eq:cr-defect-homomorphism}
\end{equation}
が成り立つ。
\end{proposition}

\begin{proof}
\(\Omega_\beta\) は部分モノイドであり、$\tau$ の両立性と反射性により \(\Gamma_\beta\) は $S\times S$ の部分モノイドである。$(a,b)\in\Gamma_\beta$ を固定する。$x\in S_\beta$ に対し、補題~\ref{cr:lem:sandwich-ideal} により $ax,bx,xa,xb$ は $S_\beta$ に属し、両立性から $ax\tau bx$ および $xa\tau xb$ である。よって系~\ref{cr:cor:local-translation} から、一意な $n_x,m_x\in N_\beta$ が存在して
\(bx=\nu_\beta(n_x)(ax)\), \(xb=\nu_\beta(m_x)(xa)\)
となる。

$x,y\in S_\beta$ に対して、結合律より $(xb)y=x(by)$。\textup{(CT)} によりその両辺はそれぞれ
\(\nu_\beta(m_x)(xay)\) と \(\nu_\beta(n_y)(xay)\)
である。作用の自由性から $m_x=n_y$。$x,y$ は独立に動くので、これらはすべて一つの共通値に一致し、これを $z_\beta(a,b)$ と書く。これで \textup{(UD)} と一意性が得られる。

乗法性を示すため、$z=z_\beta(a,b)$, $z'=z_\beta(a',b')$ とおく。$a'x\in S_\beta$ なので、
\[
\begin{aligned}
(bb')x
&=b\,\nu_\beta(z')(a'x)\\
&=\nu_\beta(z')\bigl(b(a'x)\bigr)\\
&=\nu_\beta(z')\nu_\beta(z)\bigl(a(a'x)\bigr)\\
&=\nu_\beta(zz')\bigl((aa')x\bigr).
\end{aligned}
\]
第二等号は \textup{(CE)}、最後の等号は \(N_\beta\le Z(G_\beta)\) を用いた。\textup{(UD)} における一意性から \eqref{eq:cr-defect-homomorphism} が従う。最後に $z_\beta(1,1)=1$ だから、$z_\beta$ はモノイド準同型である。
\end{proof}

\begin{corollary}[自明な欠陥]
\label{cr:cor:trivial-defect}
命題~\ref{cr:lem:uniform-defect} の仮定のもとで、
\(z_\beta(a,b)=1\)
であることと、任意の $x\in S_\beta$ に対して
\(ax=bx\) かつ \(xa=xb\)
であることは同値である。
\end{corollary}

\begin{proof}
\textup{(UD)} と中心平行移動作用の自由性から直ちに従う。
\end{proof}

\begin{lemma}[自明欠陥の局所忠実性]
\label{cr:lem:local-faithfulness}
$\tau$ が $2$-safe であり、$\beta$ が受理成分で、$a,b\in S_\beta$ が $a\tau b$ を満たすとする。もし $z_\beta(a,b)=1$ なら、$a=b$ である。
\end{lemma}

\begin{proof}
補題~\ref{cr:lem:index-rigidity} により、$a,b$ は同じ $\mathcal H$-類に属する。completely simple 半群ではこの類は群である。その単位元を $e_H$ とする。系~\ref{cr:cor:trivial-defect} から
$be_H=ae_H$、したがって $b=a$。
\end{proof}

\begin{lemma}[欠陥の乗法性]
\label{cr:lem:defect-product}
\(a_1\tau b_1,\ldots,a_m\tau b_m\) とし、\(\beta\) が表示されたすべての元の support 以下にあるとする。このとき
\[
z_\beta(a_1\cdots a_m,b_1\cdots b_m)
=\prod_{j=1}^m z_\beta(a_j,b_j).
\]
\end{lemma}

\begin{proof}
各 $(a_j,b_j)$ は \(\Gamma_\beta\) に属するので、これは $z_\beta$ の準同型性そのものである。
\end{proof}

\begin{corollary}[文脈的欠陥因子分解]
\label{cr:cor:context-defect}
\(E=u_0\blank_1u_1\cdots\blank_du_d\) を $S$ 上のタプル文脈とし、$1\le i\le d$ に対して $s_i\tau t_i$ とする。また \(\beta\) がすべての $u_j,s_i,t_i$ の support 以下にあるとする。このとき
\[
z_\beta(E[\mathbf s],E[\mathbf t])
=\prod_{i=1}^d z_\beta(s_i,t_i).
\tag{CD}
\]
\end{corollary}

\begin{proof}
\(\Gamma_\beta\) 内で
\[
(E[\mathbf s],E[\mathbf t])
=(u_0,u_0)(s_1,t_1)(u_1,u_1)\cdots(s_d,t_d)(u_d,u_d).
\]
準同型 $z_\beta$ を適用する。各固定対は系~\ref{cr:cor:trivial-defect} により欠陥 $z_\beta(u_j,u_j)=1$ をもつので、座標欠陥だけが残る。
\end{proof}

\subsection{任意タプルの二項圧縮}

\begin{lemma}[タプル文脈における support の管理]
\label{cr:lem:support-bookkeeping}
$E$ を $S$ 上のタプル文脈とする。
\(E=u_0\blank_1u_1\cdots\blank_du_d\) と書き、
\(r:=E[s_1,\ldots,s_d]\)
とおく。このとき半束 $Y$ において
\[
\operatorname{supp}(r)
=\operatorname{supp}(u_0)\operatorname{supp}(s_1)\operatorname{supp}(u_1)\cdots
\operatorname{supp}(s_d)\operatorname{supp}(u_d)
\]
である。特に、\(\operatorname{supp}(r)\) は埋め込まれた文脈に現れる各因子の support 以下にある。二つのタプルの座標 support の積が等しければ、それらを同じ文脈に埋め込んだ結果は同じ completely simple 成分に属する。
\end{lemma}

\begin{proof}
これは \textup{(CR1)} の反復適用である。成分 $S_\alpha$ と $S_\delta$ の元の積は $S_{\alpha\delta}$ に入る。$Y$ は半束なので、全因子の support の積は各因子の support 以下にある。最後の主張は、固定因子からの寄与が両方の埋め込みで共通であることから従う。
\end{proof}

次の補題が用いるのは、すでに証明した欠陥計算、補題~\ref{cr:lem:support-bookkeeping} の support 管理、および二項分離性だけである。ここから先、外部の完全正則半群定理は用いない。

\begin{lemma}[積圧縮]
\label{cr:lem:product-compression}
$\tau$ が $2$-safe であるとする。$d\ge2$ とし、
\(\mathbf s=(s_1,\ldots,s_d), \qquad \mathbf t=(t_1,\ldots,t_d)\)
が
\(s_i\tau t_i \qquad(1\le i\le d)\)
を満たすとする。二つのタプルが共通の受理タプル文脈をもつなら、
\begin{equation}
s_1s_2\cdots s_d
=
t_1t_2\cdots t_d.
\label{eq:cr-product-compression}
\end{equation}
\end{lemma}

\begin{proof}
$E$ を共通受理タプル文脈とする。両立性により \(E[\mathbf s]\tau E[\mathbf t]\)。両方の値は $P$ に属し、$2$-safe 性は unary safety を含む。関係づけられた対 \(E[\mathbf s],E[\mathbf t]\) に対して補題~\ref{cr:lem:syntactic-lifting} をアリティ1で適用し、空の単項文脈で評価すると
\begin{equation}
E[\mathbf s]
=
E[\mathbf t]
=
p\in P.
\label{eq:cr-shared-accepting-value}
\end{equation}
を得る。

\(\beta:=\operatorname{supp}(p)\) とおく。補題~\ref{cr:lem:support-bookkeeping} により、$\beta$ はすべてのタプル成分 $s_i,t_i$ および $E$ に現れるすべての固定因子の support 以下にある。

$p\in P\cap S_\beta$ であり、$\beta$ はすべての $s_i,t_i$ の support 以下にあるから、命題~\ref{cr:lem:uniform-defect} を各座標対に適用できる。\(z_i:=z_\beta(s_i,t_i)\) とおく。\eqref{eq:cr-shared-accepting-value} の二つの埋め込み積は同じ元 $p$ である。さらに $\beta=\operatorname{supp}(p)$ だから欠陥 $z_\beta(p,p)$ は定義され、命題~\ref{cr:lem:uniform-defect} の一意性から $z_\beta(p,p)=1$。したがって系~\ref{cr:cor:context-defect} より
\begin{equation}
z_1z_2\cdots z_d
=
1_{G_\beta}.
\label{eq:cr-defect-product-one}
\end{equation}
ここで積の順序は固定された Clark--Wurm のタプル順序そのものであり、文脈的因子分解の系は途中の固定因子をすでにすべて織り込んでいる。

\(A:=s_1\cdots s_{d-1}, \qquad B:=s_d, \qquad A':=t_1\cdots t_{d-1}, \qquad B':=t_d\)
とおく。両立性から
\(A\tau A', \qquad B\tau B'\)。
乗法性と \eqref{eq:cr-defect-product-one} により
\(z_\beta(AB,A'B')=1\)。
したがって系~\ref{cr:cor:trivial-defect} により、任意の $x\in S_\beta$ について
\(xAB=xA'B'\)。ゆえにすべての $x,y\in S_\beta$ に対し
\begin{equation}
xABy=xA'B'y.
\label{eq:cr-sandwich-equality}
\end{equation}
である。

さらに support の積は半束積なので、
\(\operatorname{supp}(AB)=\prod_{i=1}^d\operatorname{supp}(s_i)\)。
$\beta=\operatorname{supp}(p)$ はこの元と、共通受理文脈に現れる固定因子の support との積だから、
\(\beta\le\operatorname{supp}(AB)\)。
$c=AB$ に補題~\ref{cr:lem:sandwich-ideal} を適用すると
\(S_\beta(AB)S_\beta=S_\beta\)。
$p\in S_\beta$ なので、
\(xABy=p\)
を満たす $x,y\in S_\beta$ を選べる。すると \eqref{eq:cr-sandwich-equality} から
\(xA'B'y=p\)
でもある。したがって、座標ごとに $\tau$-関係にある二項タプル $(A,B)$ と $(A',B')$ は、文脈
$E_0=x\blank_1\blank_2y$
によって受理値 $p$ で等値化される。補題~\ref{cr:lem:equalizing-context} を適用し、固定因子なしの文脈 $\blank_1\blank_2$ で評価すると $AB=A'B'$ を得る。これはちょうど \eqref{eq:cr-product-compression} である。
\end{proof}

\begin{theorem}[compatible tolerance のアリティ2飽和]
\label{cr:thm:tolerance-saturation}
$(S,P)$ を有限の統語的点付きモノイドとし、$S$ は完全正則であるとする。$S$ 上の任意の $2$-safe compatible tolerance は、すべての有限タプル・アリティで safe である。
\end{theorem}

\begin{proof}
背理法により、$\tau$ がある有限タプル・アリティで safe でないと仮定する。$\tau$ は $2$-safe なので、この失敗はある $d\ge3$ で起こらなければならない。したがって、座標ごとに $\tau$-関係にあるタプル
\(\mathbf s=(s_1,\ldots,s_d), \qquad \mathbf t=(t_1,\ldots,t_d)\)
で、共通の受理タプル文脈をもちながらタプル分布が異なるものが存在する。

中心欠陥準同型とその乗法的帰結により補題~\ref{cr:lem:product-compression} を適用でき、
\begin{equation}
c
:=
s_1\cdots s_d
=
t_1\cdots t_d
\label{eq:cr-saturation-product}
\end{equation}
を得る。したがって support の積も一致し、
\begin{equation}
\alpha
:=
\operatorname{supp}(s_1)\cdots\operatorname{supp}(s_d)
=
\operatorname{supp}(t_1)\cdots\operatorname{supp}(t_d).
\label{eq:cr-saturation-support}
\end{equation}

必要なら二つのタプルを交換し、区別するタプル文脈 $F$ を
\begin{equation}
r:=F[\mathbf s]\in P,
\qquad
r':=F[\mathbf t]\notin P
\label{eq:cr-distinguishing-context}
\end{equation}
となるように選ぶ。$F$ に現れる固定因子の support の半束 $Y$ における積を $\kappa$ とする。固定因子がなければ、上で固定した空積の値
$\kappa=1_Y=\operatorname{supp}(1_S)$
とする。\eqref{eq:cr-saturation-support} と補題~\ref{cr:lem:support-bookkeeping} により、二つの埋め込み積は同じ completely simple 成分
\begin{equation}
S_\gamma,
\qquad
\gamma:=\kappa\alpha
\label{eq:cr-gamma-component}
\end{equation}
に属する。さらに
\(\gamma \le \operatorname{supp}(s_i), \operatorname{supp}(t_i) \qquad(1\le i\le d)\)。
$r\in P\cap S_\gamma$ なので、$\gamma$ において局所剛性補題と中心欠陥命題を適用できる。

\(z_i:=z_\gamma(s_i,t_i)\) とおく。\eqref{eq:cr-saturation-product} により、順序付きの二つの座標積は同じ元 $c$ である。さらに
\(\operatorname{supp}(c)=\alpha, \qquad \gamma=\kappa\alpha\le\alpha\)
だから $c$ は下位成分 $S_\gamma$ に作用し、欠陥 $z_\gamma(c,c)$ が定義される。命題~\ref{cr:lem:uniform-defect} の一意性により $z_\gamma(c,c)=1$。欠陥の乗法性から
\begin{equation}
z_1z_2\cdots z_d
=
1_{G_\gamma}.
\label{eq:cr-saturation-defect-one}
\end{equation}

ここでタプル文脈 $F$ を評価する。\eqref{eq:cr-gamma-component} により、$\gamma$ は $F$ の各固定因子および各タプル座標の support 以下にあり、かつ $r\in P\cap S_\gamma$。したがって系~\ref{cr:cor:context-defect} を成分 $S_\gamma$ で適用できる。共通の固定因子は欠陥 $1$ を寄与し、\eqref{eq:cr-saturation-defect-one} により
\begin{equation}
z_\gamma(r,r')
=
\prod_{i=1}^d z_i
=
1.
\label{eq:cr-context-defect-one}
\end{equation}
よって座標全体の総合的な差異は、\eqref{eq:cr-saturation-product} で用いた固定因子なしの文脈においてだけでなく、区別文脈 $F$ の固定因子を挿入した後でも自明である。

両立性から \(r\tau r'\)。両者はともに受理成分 $S_\gamma$ に属するので、\eqref{eq:cr-context-defect-one} と補題~\ref{cr:lem:local-faithfulness} から $r=r'$。これは \eqref{eq:cr-distinguishing-context} に矛盾する。

したがって $\tau$ はすべての有限タプル・アリティで safe である。
\end{proof}

\begin{theorem}[完全正則モノイドのアリティ2飽和]
\label{cr:thm:arity-two-saturation}
$L$ を正則言語とし、その統語モノイド $S$ が完全正則であるとする。このとき
\(\cmp_2(L)=\cmp_3(L)=\cmp_4(L)=\cdots.\)
さらに強く、任意の二項分離的な関係的射
$\rho:S\relto M$
は、すべての有限タプル・アリティで分離的である。
\end{theorem}

\begin{proof}
\(\rho:S\relto M\) を二項分離的とする。補題~\ref{alg:lem:relational-overlap} により、その重なり関係 \(\tau_\rho\) は $2$-safe compatible tolerance である。定理~\ref{cr:thm:tolerance-saturation} により \(\tau_\rho\) はすべての有限アリティで safe であり、同じ補題を逆向きに用いると \(\rho\) はすべての有限アリティで分離的である。$f\ge2$ に対する
\(\cmp_f(L)=\cmp_2(L)\)
は定理~\ref{alg:thm:relational-characterization} から従う。
\end{proof}

\begin{corollary}[完全正則モノイドに対する一様な高さ上界]
\label{cr:cor:height-upper}
有限完全正則モノイド全体の擬多様体 \(\mathbf{CR}\) に対して、
\(\operatorname{ch}(\mathbf{CR})\le 2\)
である。
\end{corollary}

\begin{proof}
これは定理~\ref{cr:thm:arity-two-saturation} の一様版である。
\end{proof}

\begin{remark}[なぜこの証明は orthogroup を越えて機能するのか]
冪等元に対する閉性仮定は一切用いていない。orthogroup の議論では、冪等な support の積を通してタプル成分を局所化できる。一般の完全正則モノイドでは、その近道は使えない。ここではその代わりに、下位の受理 completely simple 成分上の一様な中心平行移動欠陥を用いる。積圧縮により、一つの高アリティ共通受理を厳密な恒等式
\(s_1\cdots s_d=t_1\cdots t_d\)
へ変換し、中心欠陥計算によって、その後のあらゆるタプル文脈を同期させる。これにより completely simple 成分間の nonorthodox な相互作用が生む障害を除去できる。
\end{remark}


\section{大域的圧縮高さ三分法と鋭い領域}
\label{sec:global-arity-trichotomy}

Section~\ref{sec:vab-arity} の記法
\(
 M_{ab}=\operatorname{Synt}(\{ab\})\cong\Obj_1(\{ab\}),
 \qquad \mathbf V_{ab}=\langle M_{ab}\rangle
\)
を引き続き用いる。また
\(U=\{1,u,0\},\qquad u^2=0\)
とし、$0$ は吸収元とする。

\begin{lemma}[任意の有限非完全正則モノイドは divisor として $U$ をもつ]
\label{global:lem:U-divisor}
有限モノイド $M$ が完全正則でなければ、$U\prec M$ である。
\end{lemma}

\begin{proof}
どの部分群にも属さない $x\in M$ を選び、
\(C=\langle x\rangle^1,\qquad I=\{x^n:n\ge2\}\)
とおく。このとき $I$ は $C$ のイデアルである。また $1$ も $x$ も $I$ には属さない。
実際、ある $n\ge2$ に対して $x^n=1$ なら $x$ は単元になる。ある $n\ge2$ に対して $x^n=x$ なら、$n=2$ の場合 $x$ は冪等元であり、したがって自明な部分群 $\{x\}$ に属する。$n\ge3$ の場合、$e=x^{n-1}$ とおく。関係 $x^n=x$ から
$e^2=e$, $ex=xe=x$、さらに
$e x^{n-2}=x^{n-2}e=x^{n-2}$、
$x x^{n-2}=x^{n-2}x=e$
が従う。したがって $x^{n-2}$ は $e$ を単位元とする意味で $x$ の逆元であり、$x$ は $C$ の中で単位元 $e$ をもつ部分群に属する。いずれも $x$ の選び方に矛盾する。よって Rees 商 $C/I$ は三元
$1,\bar x,0$ をもち、$\bar x^2=0$ であるから、$U$ と同型である。
\end{proof}

\begin{lemma}[冪零性と非可換性から $M_{ab}$ が生じる]
\label{global:lem:mixing}
$M$ が有限非可換モノイドなら、$M_{ab}\prec U^2\times M$ である。
\end{lemma}

\begin{proof}
$gh\ne hg$ を満たす $g,h\in M$ を選ぶ。$U^2\times M$ において
\(\alpha:=((u,1),g)\), \(\beta:=((1,u),h)\)
とおき、
\(S:=\langle\alpha,\beta\rangle^1\)
とする。$\alpha,\beta$ の語を $U$ の第一成分へ射影すると、$\alpha$ の出現回数が $0,1,2$ 以上であるのに応じて $1,u,0$ となる。第二の $U$ 成分は $\beta$ の出現回数について同じ三分法を記録する。したがって、どちらの marker 座標にも $0$ をもたない語で表される元の出現回数は
$(0,0),(1,0),(0,1),(1,1)$
だけである。それぞれ $1$, $\alpha$, $\beta$、および $\alpha\beta$ または $\beta\alpha$ で表される。最初の三つは marker 座標によって区別される。出現回数 $(1,1)$ の二語は marker 座標は同じだが、第三座標が $gh$ と $hg$ で異なるため、$\alpha\beta\ne\beta\alpha$ である。

\(
I=S\setminus\{1,\alpha,\beta,\alpha\beta\}
\)
とおく。これが両側イデアルであることを明示的に確認する。$I\setminus\{\beta\alpha\}$ の任意の元は少なくとも一つの marker 座標に $0$ をもち、$U^2$ のある座標が一度 $0$ になれば、その後左右どちらから掛けても $0$ のままである。例外的な元 $\beta\alpha$ は各生成元を一回ずつ含む。これに $S$ の非単位元を左右いずれかから掛けると、$\alpha$ または $\beta$ のどちらかの出現回数がさらに一回増え、その marker 座標に $0$ が生じる。単位元を掛けた場合は $\beta\alpha$ 自身が $I$ に残る。したがって $SI\cup IS\subseteq I$。

よって Rees 商 $S/I$ はちょうど五元
$1$, $\bar\alpha$, $\bar\beta$, $\bar\alpha\bar\beta$, $0$
をもつ。marker の出現回数から
\(\bar\alpha^2=\bar\beta^2=0\)、また $\beta\alpha\in I$ なので
\(\bar\beta\bar\alpha=0\)。長さ3以上の任意の積は少なくとも一方の生成元を繰り返すため $0$ になる。残る非零積は $\bar\alpha\bar\beta$ だけである。これはちょうど五元の統語モノイド
\(M_{ab}=\operatorname{Synt}(\{ab\})\)
である。したがって $M_{ab}$ は部分モノイド $S\le U^2\times M$ の商であり、ゆえに
$M_{ab}\prec U^2\times M$。
\end{proof}

\begin{theorem}[\(\mathbf V_{ab}\) の構造的二分法]
\label{global:thm:structural}
\(\mathbf V\) を有限モノイドの擬多様体とする。
\(\mathbf V_{ab}\not\subseteq\mathbf V\)
ならば、
\(\mathbf V\subseteq\mathbf{Com}\)
または
\(\mathbf V\subseteq\mathbf{CR}\)
である。同値に、非可換モノイドと非完全正則モノイドをともに含む任意の擬多様体は \(\mathbf V_{ab}\) を含む。
\end{theorem}

\begin{proof}
$\mathbf V$ がどちらのクラスにも含まれないとする。非可換な $M\in\mathbf V$ と非完全正則な $R\in\mathbf V$ を選ぶ。補題~\ref{global:lem:U-divisor} と divisor に関する閉性から $U\in\mathbf V$。有限直積・部分モノイド・商に関する閉性と補題~\ref{global:lem:mixing} から $M_{ab}\in\mathbf V$、したがって
$\mathbf V_{ab}\subseteq\mathbf V$。
\end{proof}

\begin{remark}[擬多様体分類との関係]
上の divisor argument は SCL 圧縮高さに必要な障害を切り出すものであり、小さな非可換擬多様体に対する古典的分類と並ぶ位置にある。Margolis と Pin は最小非可換 varieties を分類しており、その中で
\(\operatorname{Synt}(\{ab\})\)
が生成する variety は標準的な aperiodic 例の一つである~\cite{MargolisPin1984}。Thumm の 2025 年プレプリントは
\(x_1\cdots x_n\approx\rho(x_1,\ldots,x_n)\)
という形の積恒等式を満たす擬多様体を分類している~\cite{Thumm2025}。一方、Thumm--Wei\ss の STACS 2026 の公刊論文は straight-line-program 圧縮を扱う~\cite{ThummWeiss2026}。ここで用いる分類対象はそれらとは異なり、Clark--Wurm SCL 圧縮階層の安定化スペクトルである。
\end{remark}

\begin{theorem}[大域的 SCL 圧縮高さ三分法]
\label{global:thm:trichotomy}
有限モノイドの任意の擬多様体 \(\mathbf V\) に対して
\(
\operatorname{ch}(\mathbf V)\in\{1,2,\infty\}.
\)
さらに正確には、
\[
\operatorname{ch}(\mathbf V)=\infty
\quad\Longleftrightarrow\quad
M_{ab}\in\mathbf V
\quad\Longleftrightarrow\quad
\mathbf V_{ab}\subseteq\mathbf V.
\tag{G}
\]
したがって圧縮高さが有限値 $3,4,5,\ldots$ となる擬多様体は存在しない。
\end{theorem}

\begin{proof}
まず $M_{ab}\in\mathbf V$ とする。$\mathbf V$ は擬多様体なので、これは
$\mathbf V_{ab}\subseteq\mathbf V$
と同値である。系~\ref{height:cor:core-bounds} により
\(\operatorname{ch}(\mathbf V_{ab})=\infty\)
であり、単調性から
\(\operatorname{ch}(\mathbf V)=\infty\)
が強制される。

逆に $M_{ab}\notin\mathbf V$、同値に
$\mathbf V_{ab}\not\subseteq\mathbf V$
とする。定理~\ref{global:thm:structural} により
\(
\mathbf V\subseteq\mathbf{Com}
\qquad\text{or}\qquad
\mathbf V\subseteq\mathbf{CR}.
\)
前者では定理~\ref{group:thm:commutative-collapse} により
\(\operatorname{ch}(\mathbf V)=1\)。
後者では系~\ref{cr:cor:height-upper} により
\(\operatorname{ch}(\mathbf V)\le2\)。
したがって有限値は $1$ または $2$ に限られ、無限大となるのはちょうど $M_{ab}\in\mathbf V$ のときである。
\end{proof}

\begin{remark}[三分法の射程]
この定理は、\(\operatorname{ch}\) が取りうる\emph{値}の三分法であり、同時に無限大の場合を完全に特徴づけている。ただし二つの有限の場合について完全な構造分類を主張しているわけではない。本稿では
\(
\mathbf V\subseteq\mathbf{Com}
\quad\Longrightarrow\quad
\operatorname{ch}(\mathbf V)=1
\)
および
\(
\mathbf V\subseteq\mathbf{CR}
\quad\Longrightarrow\quad
\operatorname{ch}(\mathbf V)\le2
\)
を証明するが、高さ1の逆向きの特徴づけ、したがって高さ1と2の一般的境界は未解決である。Section~\ref{sec:discussion} を参照されたい。
\end{remark}

\subsection{群統語モノイド：関係的圧縮は商へ崩壊する}

\begin{theorem}[群における商への崩壊]
\label{group:thm:quotient-collapse}
$L$ を正則言語とし、その統語モノイドが有限群 $G$ であるとする。このとき任意の $f\ge1$ に対して
\(\cmp_f(L)=\qcmp_f(L)\)。
さらに正確には、$k$ 個の値をもつ任意の有限 witness 圧縮写像から、元数が高々 $k$ の $G$ の witness 商群を取り出せる。
\end{theorem}

\begin{proof}
\(\eta:\Sigma^*\twoheadrightarrow G\)
を統語射とし、
\(h:\Sigma^*\to M\)
を任意の有限 witness 圧縮写像とする。ただし
$M=\operatorname{im}(h)$。
結合像
\(R:=\operatorname{im}(\eta,h)\le G\times M\)
を作る。第一射影
\(\pi_1:R\twoheadrightarrow G\)
は全射である。

$R$ の最小イデアルを $I$ とする。像 $\pi_1(I)$ は群 $G$ の空でないイデアルなので、$\pi_1(I)=G$。冪等元 $e\in I$ を選ぶ。群の冪等元は一つしかないから、$\pi_1(e)=1_G$。有限半群の最小イデアルの標準構造により、
\(H:=eIe\)
は単位元 $e$ をもつ群である~\cite{Howie1995}。さらに $\pi_1(H)=G$。実際、$g\in G$ に対して $\pi_1(x)=g$ を満たす $x\in I$ を選べば、$exe\in H$ かつ $\pi_1(exe)=g$ である。

\(
\psi:=\pi_2|_H,
\qquad H_M:=\psi(H),
\qquad K:=\ker\psi
\)
とおく。このとき $H_M$ は $M$ に含まれる有限群である。$\pi_1|_H:H\twoheadrightarrow G$ は全射なので、
\(N:=\pi_1(K)\trianglelefteq G\)。
さらに全射準同型
\(
H_M\twoheadrightarrow G/N,
\qquad
\psi(x)\longmapsto \pi_1(x)N
\)
が well-defined に存在する。実際、$\psi(x)=\psi(y)$ なら $x^{-1}y\in K$ なので、
$\pi_1(x)^{-1}\pi_1(y)\in N$。
したがって
\(|G/N|\le |H_M|\le |M|\)。
残るのは $G\to G/N$ がすべての unsafe tuple を分離することを示すことである。逆に、ある
$(\mathbf s,\mathbf t)\in\mathcal U_d(G,P)$, $d\le f$
が各座標で潰される、すなわちすべての $i$ に対して $s_i^{-1}t_i\in N$ と仮定する。各 $i$ について $\pi_1(k_i)=s_i^{-1}t_i$ を満たす $k_i\in K$ を選び、また $\pi_1(a_i)=s_i$ を満たす $a_i\in H$ を選ぶ。すると
\(
\pi_1(a_ik_i)=t_i,
\qquad
\psi(a_ik_i)=\psi(a_i).
\)
$H\subseteq R=\operatorname{im}(\eta,h)$ なので、共通の第二座標 $\psi(a_i)$ は
\(\rho_h(s_i)\cap\rho_h(t_i)\)
に属する。これがすべての座標で起こるが、定理~\ref{alg:thm:safety-interface} により $\rho_h$ は $f$-分離的であるから矛盾する。

したがって $G/N$ は大きさ高々 $|M|$ の witness 商である。よって
\(\qcmp_f(L)\le\cmp_f(L)\)
であり、逆向きの不等式は定義から成り立つ。
\end{proof}

\subsection{高さ1を与える第二の機構：バンド}

可換性は unary collapse に必要ではない。有限バンド、すなわち高度に非可換なものも含む冪等モノイド全体の擬多様体は、すでにアリティ1で飽和する。標準的な、バンドを長方形 $\mathcal D$-類の半束へ分解する表現を用いる。

\begin{theorem}[バンドの単項飽和]
\label{band:thm:unary-saturation}
$L$ を正則言語とし、その統語モノイド $S$ が有限バンドであるとする。このとき $S$ からの任意の単項分離的な関係的射は、すべての有限アリティで分離的である。したがって
\(\cmp_1(L)=\cmp_2(L)=\cmp_3(L)=\cdots.\)
\end{theorem}

\begin{proof}
バンドの標準的な半束分解を
\(S=\bigsqcup_{\alpha\in Y}S_\alpha\)
と書く。各 $S_\alpha$ は長方形バンドであり、
\(S_\alpha S_\beta\subseteq S_{\alpha\beta}\)
である。まず証明を駆動する基本的な吸収則を記す。
$a\in S_\alpha$, $p\in S_\beta$, $\beta\le\alpha$ とすると、$pap\in S_\beta$。$S_\beta$ は長方形バンドなので
$p(pap)p=p$。
一方、結合律と $p$ の冪等性により左辺は $pap$ に等しい。したがって
\[
 pap=p.
\tag{B}
\]

\(\rho:S\relto M\)
を単項分離的とし、
\(\rho(x)\cap\rho(y)\ne\varnothing\)
のとき $x\tau y$ と書く。座標ごとに $\tau$-関係にある二つの $d$-タプルが共通の受理タプル文脈をもつと仮定する。両立性により、二つの受理済み埋め込み値は $\tau$-関係にある。単項分離性と統語的識別可能性から、その二つの値は同一の元 $p\in P$ でなければならない。
\(\beta=\operatorname{supp}(p)\)
とおく。半束における support の積から、$\beta$ は各タプル座標の support 以下にある。したがって \textup{(B)} により各座標対 $s_i\tau t_i$ に対して
\(p s_i p=p=p t_i p\)。
よって $s_i$ と $t_i$ は共通の受理単項文脈
$p\blank_1p$
をもつ。もし両者が異なれば、統語性によりそれらを区別する別の単項文脈が存在するから、$\rho$ によって潰される unsafe な単項対を形成し、単項分離性に反する。したがってすべての $i$ で $s_i=t_i$。二つのタプルは同一であり、どのアリティでも unsafe になりえない。ゆえに $\rho$ はすべての有限アリティで分離的である。
\end{proof}

\subsection{核群言語：アリティ2は内部共役を検出する}

\(\eta:\Sigma^*\twoheadrightarrow G\)
を全射とし、
\(K_\eta:=\eta^{-1}(\{1_G\})\)
とおく。

\begin{theorem}[中心指数定理]
\label{group:thm:center-index}
$K_\eta$ の統語モノイドは $G$ であり、
\(
\cmp_1(K_\eta)=1,
\qquad
\cmp_f(K_\eta)=[G:Z(G)]=|\operatorname{Inn}(G)|
\quad(f\ge2).
\)
\end{theorem}

\begin{proof}
異なる群元は統語的に区別可能である。$g\neq h$ とし、$\eta(w_g)=g^{-1}$ を満たす語 $w_g\in\Sigma^*$ を選ぶ。統語値 $g$ をもつ任意の語に $w_g$ を後置すると受理されるが、統語値 $h$ をもつ語では拒否される。したがって統語モノイドは $G$ である。

アリティ1では、$s,t\in G$ が共通の受理文脈をもつなら、ある $p,q\in G$ に対して
$psq=ptq=1_G$。消去則から $s=t$。したがって異なる unsafe な単項対は存在せず、
\(\cmp_1(K_\eta)=1\)。

任意の $f\ge2$ に対する上界には商
$G\twoheadrightarrow G/Z(G)$
を用いる。$t_i=s_i z_i$ かつ $z_i\in Z(G)$ とする。任意のタプル文脈において、すべての $z_i$ は固定因子およびタプル成分を越えて移動できるので、第二の埋め込みの群値は第一の埋め込み値に $z_1\cdots z_d$ を掛けたものになる。ある一つの文脈が両タプルを受理するなら
$z_1\cdots z_d=1_G$。
したがって両者は任意のタプル文脈で同じ群値をもつ。よって
\(\cmp_f(K_\eta)\le [G:Z(G)]\)。

下界については、定理~\ref{group:thm:quotient-collapse} により、アリティ2を witness する商
$G\to G/N$
だけを考えればよい。もし $g\in N\setminus Z(G)$ なら、$gk\neq kg$ を満たす $k\in G$ を選び、
\(
\eta(u)=k,
\qquad
\eta(v)=k^{-1}
\)
を満たす語 $u,v\in\Sigma^*$ を選ぶ。タプル
\((1_G,1_G), \qquad (g,g^{-1})\)
は各座標で $N$ を法として等しく、隣接した穴をもつ受理文脈を共有する。しかしタプル文脈
\(F=\blank_1\,u\,\blank_2\,v\)
は第一タプル上で統語値 $1_G$、第二タプル上で
$gkg^{-1}k^{-1}\neq1_G$
を与える。したがって両タプルは unsafe であり、矛盾する。ゆえに
$N\le Z(G)$
でなければならず、任意の witness 商の大きさは少なくとも
$[G:Z(G)]$。以上から等号が従う。
\end{proof}

\begin{corollary}[高さ1境界の縮約]
\label{height:cor:height-one-reduction}
\(\mathbf V\) を有限モノイドの擬多様体とする。
\begin{enumerate}[label=\textup{(\roman*)},leftmargin=*]
\item \(\operatorname{ch}(\mathbf V)<\infty\) なら、
\(
 \mathbf V\subseteq\mathbf{Com}
 \text{ or }
 \mathbf V\subseteq\mathbf{CR}.
\)
\item \(\mathbf V\subseteq\mathbf{CR}\) かつ \(\mathbf V\) が非可換有限群を含むなら、
\(\operatorname{ch}(\mathbf V)=2\)。
\end{enumerate}
したがって、すでに解決済みの可換な場合を除けば、高さ1分類の未解決部分は、群として現れる元がすべて可換群であるような擬多様体
\(\mathbf V\subseteq\mathbf{CR}\)
に限られる。
\end{corollary}

\begin{proof}
\textup{(i)} について、有限高さと定理~\ref{global:thm:trichotomy} から
$M_{ab}\notin\mathbf V$。よって定理~\ref{global:thm:structural} が主張の二分法を与える。
\textup{(ii)} について、系~\ref{cr:cor:height-upper} により
\(\operatorname{ch}(\mathbf V)\le2\)。
非可換有限群 $G\in\mathbf V$ をとり、有限アルファベット上の全射
\(\eta:\Sigma^*\twoheadrightarrow G\)
を選ぶ。核言語 $K_\eta$ の統語モノイドは $G$ であり、定理~\ref{group:thm:center-index} により
\(
 \cmp_1(K_\eta)=1< [G:Z(G)]=\cmp_2(K_\eta),
\)
なぜなら非可換群の中心は真部分群だからである。したがって高さ1は不可能であり、上界は鋭い。
\end{proof}

\subsection{高さの鋭さと古典的ベンチマーク}

\begin{corollary}[バンドの擬多様体は高さ1をもつ]
\label{band:cor:height-one}
有限バンド全体の擬多様体を \(\mathbf B\) とする。このとき
\(\operatorname{ch}(\mathbf B)=1\)。
特に、高さ1は可換性を意味しない。標準的な最小非可換 one-sided band 擬多様体
\(\langle U_2\rangle\) および \(\langle U_2^{\mathrm{op}}\rangle\)
は \(\mathbf B\) に含まれ、したがってこれらも高さ1をもつ。
\end{corollary}

\begin{proof}
高さの主張は定理~\ref{band:thm:unary-saturation} である。三元モノイド
\(U_2=\operatorname{Synt}(A^*a)\)
は非可換バンドであり、その反対モノイドは $aA^*$ の統語モノイドである~\cite{MargolisPin1984}。
\end{proof}

\begin{corollary}[完全正則上界の鋭さ]
\label{cr:cor:height}
有限完全正則モノイド全体の擬多様体 $\mathbf{CR}$ に対して、
\(
\operatorname{ch}(\mathbf{CR})=2.
\)
\end{corollary}

\begin{proof}
定理~\ref{cr:thm:arity-two-saturation} が上界を与える。擬多様体 $\mathbf{CR}$ は非可換有限群を含む。例えば
$Z(S_3)=1$
なので、統語群 $S_3$ に対する核言語は中心指数定理により
\(\cmp_1=1, \qquad \cmp_f=6\quad(f\ge2)\)
を満たす。したがってアリティ1からアリティ2への真の跳躍が起こり、高さはちょうど2である。
\end{proof}

\subsection{ベンチマーク擬多様体と最小障害}

\begin{corollary}[ベンチマーク擬多様体]
\label{global:cor:benchmarks}
有限群全体を \(\mathbf G\)、有限 aperiodic モノイド全体を \(\mathbf A\)、
$\mathcal J$-, $\mathcal R$-, $\mathcal L$-trivial 擬多様体をそれぞれ
\(\mathbf J,\mathbf R,\mathbf L\)
とする標準記法のもとで、次の値が成り立つ：
\[
\begin{array}{c|cccccccc}
\mathbf V
&\mathbf{Com}&\mathbf B&\mathbf G&\mathbf{CR}
&\mathbf A&\mathbf J&\mathbf R&\mathbf L\\
\hline
\operatorname{ch}(\mathbf V)
&1&1&2&2&\infty&\infty&\infty&\infty
\end{array}
\]
\end{corollary}

\begin{proof}
最初の値は可換崩壊から、第二の値は系~\ref{band:cor:height-one} から従う。群および完全正則モノイドの高さは高々2であり、一方、任意の非可換有限群の核言語は定理~\ref{group:thm:center-index} により
\(\cmp_1=1<\cmp_2\)
を満たすから、$\mathbf G$ と $\mathbf{CR}$ の高さはいずれもちょうど2である。

モノイド $M_{ab}$ は aperiodic かつ $\mathcal J$-trivial である（その主両側イデアルはすべて相異なる）。したがって
\(M_{ab}\in\mathbf A\cap\mathbf J\)
であり、
\(\mathbf J\subseteq\mathbf R\cap\mathbf L\)。
よって無限大の各値は定理~\ref{global:thm:trichotomy} から従う。
\end{proof}

このベンチマーク表は、古典的な Margolis--Pin の最小非可換フロンティア~\cite{MargolisPin1984} の位置づけも与える。二つの one-sided band 型は高さ1、非可換群型は高さ2、そして
\(\langle M_{ab}\rangle\)
は高さ無限大をもつ。したがって三つの値はすべて、この最小非可換フロンティア上ですでに実現される。


\section{グラフ実現と計算量}
\label{sec:graph-universality}

アリティ1の枠組みは文脈衝突グラフ \(\Gamma_L\) から始まった。ここでは、このグラフ幾何が普遍的であることを示す。すなわち、任意の空でない有限単純グラフは、長さ3の有限言語において、指定した頂点記号上の誘導部分グラフとして実現され、他のすべての衝突グラフ頂点は孤立する。したがって、そこから得られる彩色数との等式は、まず表現定理であり、その次に NP-hardness の還元を与える。

\subsection{グラフ言語の構成}

\(G=(V,E)\) を有限単純グラフとする。

各頂点 \(v\in V\) に対して終端記号 \(x_v\) を導入する。各辺 \(e\in E\) に対して二つの終端記号 \(\ell_e,r_e\) を導入する。さらに各頂点 \(v\in V\) に対して二つの private terminal \(p_v,q_v\) を導入する。これらの終端記号はすべて互いに異なるものとする。

辺 \(e=\{u,v\}\) に対して、二つの辺語
\(\ell_e x_u r_e, \qquad \ell_e x_v r_e\)
を言語に入れる。また各頂点 \(v\) に対して private word \(p_vx_vq_v\) も言語に入れる。

したがって
\(
L_G
:=
\{
\ell_ex_ur_e,\ell_ex_vr_e:
e=\{u,v\}\in E
\}
\cup
\{
p_vx_vq_v:v\in V
\}.
\)
$L_G$ のすべての語の長さはちょうど3であり、
\(|L_G|=2|E|+|V|\)。

\subsection{unsafe 因子の分類}

\begin{lemma}[共通単項文脈の分類]
\label{graph:lem:factor-classification}
二つの因子 $s,t$ が $L_G$ に対して共通の受理単項文脈をもつなら、
$|s|=|t|\le3$ である。さらに：
\begin{enumerate}[label=\textup{(\roman*)},leftmargin=*]
\item 長さ $0$ と $3$ からは相異なる unsafe pair は生じない；
\item 共通受理文脈をもつ相異なる長さ2因子は、同じ単項分布をもつ；
\item 長さ1の分布は
\[
\D_{L_G}(x_v)=\{(\ell_e,r_e):e\ni v\}\cup\{(p_v,q_v)\},
\]
\[
\D_{L_G}(\ell_e)=\{(\varepsilon,x_ur_e),(\varepsilon,x_vr_e)\},
\qquad
\D_{L_G}(r_e)=\{(\ell_ex_u,\varepsilon),(\ell_ex_v,\varepsilon)\}
\]
であり、ここで $e=\{u,v\}$。また
\(
\D_{L_G}(p_v)=\{(\varepsilon,x_vq_v)\}
\)
および
\(
\D_{L_G}(q_v)=\{(p_vx_v,\varepsilon)\}.
\)
\end{enumerate}
\end{lemma}

\begin{proof}
$psq,ptq\in L_G$ なら、完成した二語はいずれも長さ3なので、
$|s|=|t|\le3$ であり、$s,t$ はともに受理語の因子である。長さ0では $\varepsilon$ しか存在せず、各受理長さ3語の分布は
$\{(\varepsilon,\varepsilon)\}$
なので \textup{(i)} が従う。

\textup{(ii)} について、private でない共通の長さ2因子対は、$e=\{u,v\}$ に対して
\[
\D_{L_G}(\ell_ex_u)=\D_{L_G}(\ell_ex_v)=\{(\varepsilon,r_e)\},
\qquad
\D_{L_G}(x_ur_e)=\D_{L_G}(x_vr_e)=\{(\ell_e,\varepsilon)\}
\]
だけである。private 因子 $p_vx_v$ と $x_vq_v$ は、それぞれ singleton 分布
$\{(\varepsilon,q_v)\}$ と $\{(p_v,\varepsilon)\}$
をもつ。異なる gadget は、それぞれ固有の completing marker を共有できず、また穴の位置は固定されているので prefix 因子と suffix 因子が同じ文脈を共有することもない。したがって \textup{(ii)} が従う。最後に \textup{(iii)} は、$L_G$ の定義語における各終端記号の出現を列挙すれば得られる。
\end{proof}

\begin{lemma}[unsafe-pair グラフは正確に \(G\) である]
\label{graph:lem:unsafe-exact}
\(L_G\) の相異なる unsafe pair は正確に
\(\{x_u,x_v\} \quad\text{with}\quad \{u,v\}\in E\)
である。より正確には、相異なる頂点 \(u,v\) に対して
\(x_u,x_v\text{ are unsafe} \iff \{u,v\}\in E\)
であり、頂点でない終端記号を含む対は一つも unsafe ではない。
\end{lemma}

\begin{proof}
まず \(e=\{u,v\}\in E\) とする。このとき文脈
\(\ell_e\blank_1r_e\)
は $x_u$ と $x_v$ の両方を受理するので、両者の分布は交わる。一方で
\((p_u,q_u)\in\D_{L_G}(x_u)\)
であり、marker の private 性により
\((p_u,q_u)\notin\D_{L_G}(x_v)\)。
したがって分布は異なり、$x_u,x_v$ は unsafe である。

逆に $u\neq v$ かつ
\(\D_{L_G}(x_u)\cap\D_{L_G}(x_v)\neq\varnothing\)
とする。補題~\ref{graph:lem:factor-classification}\textup{(iii)} により、共通文脈は $u,v$ の両方に接続する辺 $e$ に対する
$(\ell_e,r_e)$
でなければならない。グラフは単純なので $e=\{u,v\}$。残るのは marker 終端記号を含む対を排除することである。各 $x_v$ は受理語の中央位置にしか現れない。終端記号 $\ell_e,p_v$ は第一位置にしか現れず、$r_e,q_v$ は第三位置にしか現れない。共通単項文脈は穴の位置を固定するので、どの $x_v$ も marker 終端記号と共通受理文脈をもつことはできない。

二つの第一位置 marker が共通受理文脈をもつのは、両者の後ろに同じ長さ2 suffix が続く場合だけである。$\ell_e$ の後ろの suffix は private edge marker $r_e$ で終わり、$p_v$ の後ろの suffix は private marker $q_v$ で終わる。異なる gadget に属する marker は互いに異なるので、相異なる第一位置 marker が共通受理文脈をもつことはない。第三位置 marker についても、その長さ2 prefix を用いれば同じ議論が成り立つ。

補題~\ref{graph:lem:factor-classification}\textup{(i)--(ii)} と合わせると、これですべての因子長の場合を尽くしている。
\end{proof}

\begin{corollary}[任意の有限グラフは誘導された文脈衝突幾何として現れる]
\label{graph:cor:conflict-realization}
言語 \(L_G\) に対して、指定した頂点終端記号
\(\{x_v:v\in V(G)\}\)
が \(\Gamma_{L_G}\) に誘導する部分グラフは \(G\) と同型であり、\(\Gamma_{L_G}\) の他のすべての頂点は孤立している。同値に、\(\Gamma_{L_G}\) はこの $G$ の誘導コピーと孤立頂点からなる独立集合との非交和である。
\end{corollary}

\begin{proof}
補題~\ref{graph:lem:unsafe-exact} は、受理語の因子間にあるすべての unsafe pair を同定している。受理語の因子でない語は空の分布をもち、したがってどの unsafe pair にも参加しない。主張は Section~\ref{sec:conflict} における \(\Gamma_L\) の定義から従う。

さらに、指定語 $x_v$ は互いに異なる統語類に属する。なぜなら private 文脈 $(p_v,q_v)$ は
\(\D_{L_G}(x_v)\)
には属するが、$u\ne v$ のとき
\(\D_{L_G}(x_u)\)
には属さないからである。したがって系~\ref{conf:cor:omega-syntactic} により、$G$ は単項 principal-overlap graph \(\Om_1(L_G)\) の誘導部分グラフとしても実現される。
\end{proof}

\subsection{彩色数の厳密実現}

この構成の重要性は、単に hardness が従うことではなく、彩色数との厳密な等式にある。特に、指定した長さ1因子の圧縮類はちょうどグラフの色類に対応する。残りの因子は、新たな unsafe 同一視を一切生まないように配置されている。アルファベットの大きさはグラフとともに増える。対応する固定二元アルファベット上の問題は未解決として残す。

\begin{lemma}[任意の圧縮写像は proper coloring を誘導する]
\label{graph:lem:lower-coloring}
\(h:\Sigma_G^*\to M\)
が \(L_G\) の unary substitutability を witness するなら、
\(c_h:V\to\operatorname{im}(h), \qquad c_h(v):=h(x_v)\)
は $G$ の proper coloring である。したがって
\(
\chi(G)\le|\operatorname{im}(h)|.
\)
\end{lemma}

\begin{proof}
各辺 \(\{u,v\}\in E\) に対して、補題~\ref{graph:lem:unsafe-exact} により $x_u,x_v$ は unsafe である。witness 圧縮写像は各 unsafe pair に異なる値を割り当てなければならないので、
$h(x_u)\neq h(x_v)$。したがって $c_h$ は proper である。
\end{proof}

\begin{lemma}[任意の proper coloring は圧縮写像を与える]
\label{graph:lem:upper-coloring}
$G$ が $q\ge2$ 色の proper coloring をもつなら、$L_G$ の unary substitutability を witness する $q$ 元モノイド圧縮写像が存在する。$G$ が辺をもたないなら、一元圧縮が witness となる。
\end{lemma}

\begin{proof}
まず $q\ge2$ とする。
\(M_q=\{1,0,c_1,\ldots,c_{q-2}\}\)
とおく。$1$ を単位元、$0$ を吸収元とし、$m,n\neq1$ のとき常に $mn=0$ と定める。これはちょうど $q$ 元をもつ結合的モノイドである。

$q$ 個のモノイド元を色名として使う proper coloring
\(\gamma:V\to M_q\)
を選ぶ。生成元上の圧縮写像を
\(h(x_v):=\gamma(v)\)
とし、すべての marker 終端記号 $\ell_e,r_e,p_v,q_v$ を $0$ に送る。$\Sigma_G^*$ の自由性により、これは一意にモノイド準同型
\(h:\Sigma_G^*\to M_q\)
へ延長される。補題~\ref{graph:lem:unsafe-exact} により、各 unsafe pair は辺 $\{u,v\}$ に対応する $x_u,x_v$ の形である。$\gamma$ の proper 性により $h(x_u)\neq h(x_v)$。したがって unsafe pair は一つも同じ圧縮値を受けず、$h$ は unary substitutability を witness する。

$G$ が edgeless なら、補題~\ref{graph:lem:unsafe-exact} により $L_G$ には unsafe pair が全く存在しない。したがって一元モノイドへの唯一の準同型が witness となる。
\end{proof}

\begin{theorem}[グラフの厳密実現]
\label{graph:thm:graph-realization}
任意の空でない有限単純グラフ \(G\) に対して、
\(\cmp_1(L_G)=\chi(G)\)。
さらに構成 \(G\mapsto L_G\) は多項式時間であり、$L_G$ のすべての語は長さ3、かつ
\(|L_G|=2|E|+|V|\)
である。
\end{theorem}

\begin{proof}
補題~\ref{graph:lem:lower-coloring} から
\(\chi(G)\le\cmp_1(L_G)\)。
最小 proper coloring と補題~\ref{graph:lem:upper-coloring} から
\(\cmp_1(L_G)\le\chi(G)\)。
構成サイズと語長に関する主張は $L_G$ の定義から直ちに従う。
\end{proof}

\begin{remark}[幾何学的下界の鋭さ]
定理~\ref{graph:thm:graph-realization} は、命題~\ref{scl:prop:chromatic-lower} の一般的な幾何学的下界が、すでに単項アリティで鋭いことを示す。実際、標準的な単項同一視のもとで
\(\chi(\Om_1(L_G))=\chi(G)=\cmp_1(L_G)\)。
したがって任意の有限衝突幾何は、合成的圧縮の下界を厳密に達成しうる。
\end{remark}

\begin{remark}
アルファベット \(\Sigma_G\) は、明示的な有限言語入力の一部である。固定二元アルファベット上の対応する variable-threshold 問題は未解決である。厳密実現定理は、明示的有限言語に対する自然な可変アルファベット表現を扱う。
\end{remark}

\subsection{有限言語に対する決定問題}

\begin{definition}
\({\normalfont\textsc{Length-3-Obs}}_1(3)\)
を次の決定問題とする。

\medskip
\noindent
\textbf{入力：} 有限アルファベット \(\Sigma\) も入力の一部として、明示的に列挙された有限言語
\(L\subseteq\Sigma^3\)。

\noindent
\textbf{問い：} \(\cmp_1(L)\le3\) か？
\end{definition}

\begin{lemma}[固定閾値に対する多項式時間検証]
\label{graph:lem:np}
問題
\(
{\normalfont\textsc{Length-3-Obs}}_1(3)
\)
は NP に属する。
\end{lemma}

\begin{proof}
証明書は、元数高々3のモノイド $M$ の積表、その単位元、および各終端記号の値 $h(a)\in M$ を与える。表の大きさは定数であり、モノイド則は直接検査できる。また生成元への割当ては準同型
\(h:\Sigma^*\to M\)
を決定する。

非空の単項分布をもつのは列挙された語の因子だけである。すべての受理語は長さ3なので、そのような因子は多項式個しかない。各因子 $x$ に対して完全分布
\(\D_L(x)=\{(p,q):pxq\in L\}\)
を列挙する。次にすべての因子対 $x,y$ を検査し、$h(x)=h(y)$ かつ二つの分布が交わるなら、それらが等しいことを要求する。これは正確に unary substitutability の条件であり、明示入力サイズに関して多項式時間で検査できる。
\end{proof}

\begin{remark}[NP の議論を固定閾値で述べる理由]
上の証明書による議論では、閾値が固定されていることを用いている。$k=3$ のとき witness モノイドの積表は定数サイズである。もし $k$ 自身が入力の一部なら、明示的な積表は
\(\Theta(k^2\log k)\)
bits を必要とする。言語の符号化サイズから関係する witness サイズへの事前の多項式上界がない限り、同じ議論だけでは variable-threshold 問題が NP に属するとは言えない。この点は後述の regular-input complexity 問題に含める。
\end{remark}

\begin{theorem}[明示列挙された長さ3言語に対する NP 完全性]
\label{graph:thm:npcomplete}
有限アルファベットが入力の一部であるとき、問題
\({\normalfont\textsc{Length-3-Obs}}_1(3)\)
は NP-complete である。
\end{theorem}

\begin{proof}
NP に属することは補題~\ref{graph:lem:np} による。

NP-hardness については {
ormalfont\textsc{3-Colorability}} から還元する~\cite{GareyJohnson1979}。有限単純グラフ $G$ が与えられたとき、もし $V(G)=\varnothing$ なら孤立頂点を一つ追加する。これは3彩色可能性を保存する。したがって $G$ は空でないとしてよく、上のように $L_G$ を構成する。この構成は多項式時間であり、
\(L_G\subseteq\Sigma_G^3\)
を与える。定理~\ref{graph:thm:graph-realization} により
\(\cmp_1(L_G)=\chi(G)\)。
したがって
\(
G\text{ が3彩色可能}
\iff \chi(G)\le3
\iff \cmp_1(L_G)\le3.
\)
よって還元は正しい。
\end{proof}

\subsection{明示的な点付きモノイドからの計算量}

上の hardness は、入力表現として列挙有限言語を用いたことによる人工的なものではない。有限代数上の最適化それ自体がすでに NP-complete である。

\begin{definition}
\({\normalfont\textsc{Pointed-Obs}}_1\)
を次の決定問題とする。

\medskip
\noindent
\textbf{入力：} 積表で与えられた有限モノイド \(T\)、受理部分集合 \(P\subseteq T\)、および正整数 \(k\)。

\noindent
\textbf{問い：} ある有限モノイド \(M\), \(|M|\le k\) への単項分離的関係的射
\(\rho:T\relto M\)
が存在するか？
\end{definition}

\begin{theorem}[点付きモノイド入力に対する NP 完全性]
\label{graph:thm:pointed-npcomplete}
問題 \({\normalfont\textsc{Pointed-Obs}}_1\) は NP-complete である。NP-hardness は固定閾値 $k=3$ ですでに成立し、入力が有限長さ3言語の統語的点付きモノイドであると約束されている場合にも成立する。
\end{theorem}

\begin{proof}
NP membership を示す。$n=|T|$ とおく。$k\ge n$ なら恒等射が witness なので、$k<n$ と仮定する。証明書は、元数高々 $k$ のモノイド $M$ の積表、その単位元、および関係
\(\rho\subseteq T\times M\)
からなる。$k<n$ なので、そのサイズは入力に関して多項式である。モノイド則、各 fiber の非空性、単位元条件、および包含
\(\rho(s)\rho(t)\subseteq\rho(st)\)
は多項式時間で検査できる。各 $s\in T$ に対して
\(\Phi_1(s)=\{(q_0,q_1)\in T^2:q_0sq_1\in P\}\)
を列挙する。これにより unary unsafe pair をすべて多項式時間で同定でき、その後、各 unsafe pair に対する条件
\(\rho(s)\cap\rho(t)=\varnothing\)
も多項式時間で検査できる。したがって問題は NP に属する。

NP-hardness について、有限単純グラフ $G$ から始める。$V(G)=\varnothing$ なら孤立頂点を一つ追加し、3彩色可能性を保存する。よって $G$ は空でないとして、定理~\ref{graph:thm:graph-realization} の長さ3言語 $L_G$ を構成する。その点付き統語モノイドは多項式時間で構成できる。実際、非空の単項分布をもつ語はすべて列挙長さ3語の連続因子であり、そのような因子は多項式個しかない。一方、因子でないすべての語は空の分布をもち、一つの零類をなす。統語類の等しさは明示的に列挙できる単項分布の等しさであり、類の積は代表元を連結し、その結果が属する因子類、または零類を参照することで得られる。したがって完全な積表と受理部分集合は多項式時間で計算でき、サイズも多項式である。

定理~\ref{alg:thm:relational-characterization} と統語的不変性により、この点付きモノイドからの単項分離的関係的射の最小終域サイズは正確に
\(\cmp_1(L_G)=\chi(G)\)
である。したがって $G$ が3彩色可能であることと
\({\normalfont\textsc{Pointed-Obs}}_1(T_{L_G},P_G,3)\)
が yes-instance であることは同値である。還元は多項式時間なので、NP-hardness が従う。
\end{proof}


\section{議論と未解決問題}
\label{sec:discussion}

中心的な現象は、局所的な豊かさと大域的な剛性との対比である。\(\mathbf V_{ab}\) の内部では、個々の言語が任意の有限の真アリティ初期部分と、非有界な連続ギャップを実現する。しかし擬多様体全体で一様化すると、
\(
\operatorname{ch}(\mathbf V)\in\{1,2,\infty\},
\qquad
\operatorname{ch}(\mathbf V)=\infty \Longleftrightarrow M_{ab}\in\mathbf V.
\)
完全正則飽和定理は有限側を説明する。すなわち、点付き二項分離は差異を乗法的な中心欠陥へ還元し、共通の高アリティ受理はそれらの総合欠陥を消滅させる。

SCL 定式化はまた、幾何学的複雑さと、再利用可能な最小合成資源の大きさとを分離する。非可換群の核言語では principal tuple geometry が急速に成長しうる一方、最適値はアリティ2で安定化する。単項アリティでは、グラフ実現により任意の有限衝突幾何を
\(\cmp_1(L_G)=\chi(G)\)
とともに実現できる。さらに独立に、厳密な unary の例は、圧縮を統語商を通して因子化するよう強制するとコストが非有界になりうることを示す。したがってこの三分法が扱うのは、有限合成分離の安定化であり、概念束そのものや商複雑性の安定化ではない。

特に自然と思われる問題を二つ挙げる。
\begin{enumerate}[leftmargin=*]
\item \textbf{\(\mathbf{CR}\) 内に残る高さ1の境界。}
系~\ref{height:cor:height-one-reduction} により、分類の未解決部分は、群として現れる元がすべて可換群である完全正則擬多様体へ縮約される。次は真だろうか：
\(
\mathbf V\subseteq\mathbf{CR}
\quad\text{かつ}\quad
\mathbf V\text{ のすべての群が可換}
\quad\Longrightarrow\quad
\operatorname{ch}(\mathbf V)=1\,?
\)
バンドはすべての極大部分群が自明である極端な場合を与え、可換モノイドも別の肯定的領域を与える。具体的な最初の検証対象として、有限可換群 \(A\) 上の completely simple Rees matrix semigroup
\(\mathcal M[A;I,\Lambda;Q]\)
に単位元を付加して得られる有限モノイド、およびそれらが生成する擬多様体が挙げられる。

\item \textbf{点付きモノイドモデルを越える regular-input complexity。}
定理~\ref{graph:thm:pointed-npcomplete} は、明示的点付きモノイドに対する単項問題を解決した。DFA から与えた場合の計算量、アリティ境界そのものが変数である場合の挙動、および明示的な統語モノイドの blow-up を回避できるかどうかを決定せよ。

\end{enumerate}

\section*{声明および宣言}

\noindent\textbf{著者の貢献.}
著者が研究を着想し、数学的結果を構築し、証明を検証し、原稿を作成した。

\noindent\textbf{研究資金.}
本研究に対する研究資金の提供は受けていない。

\noindent\textbf{利益相反.}
著者は利益相反がないことを宣言する。

\noindent\textbf{データ利用可能性.}
本研究は理論研究であり、データセットの生成または解析は行っていない。

\noindent\textbf{生成AIの利用.}
原稿作成の過程で、著者は文献検索、証明監査、構成整理、および言語編集の補助として OpenAI ChatGPT を使用した。数学的主張、証明、引用、および最終原稿についての責任はすべて著者に帰属する。

\small
\begin{thebibliography}{99}

\bibitem{ClarkSCL}
A.~Clark,
\newblock The syntactic concept lattice: another algebraic theory of the context-free languages?,
\newblock \emph{Journal of Logic and Computation} 25(5) (2015), 1203--1229,
\newblock \url{https://doi.org/10.1093/logcom/ext037}.

\bibitem{WurmSCLTuples}
C.~Wurm,
\newblock On some extensions of syntactic concept lattices: completeness and finiteness results,
\newblock in: \emph{Formal Grammar 2015--2016}, LNCS 9804, Springer (2016), 164--179,
\newblock \url{https://doi.org/10.1007/978-3-662-53042-9_10}.

\bibitem{WurmJLLI2017}
C.~Wurm,
\newblock Language-theoretic and finite relation models for the (Full) Lambek calculus,
\newblock \emph{Journal of Logic, Language and Information} 26(2) (2017), 179--214.
\newblock \url{https://doi.org/10.1007/s10849-017-9249-z}.

\bibitem{Kuznetsov2026}
S.~L.~Kuznetsov,
\newblock On syntactic concept lattice models for the Lambek calculus and infinitary action logic,
\newblock \emph{Journal of Logic and Computation} 36(1) (2026), Article exaf078.
\newblock \url{https://doi.org/10.1093/logcom/exaf078}.

\bibitem{ClarkYoshinaka2014}
A.~Clark and R.~Yoshinaka,
\newblock An algebraic approach to multiple context-free grammars,
\newblock in: \emph{Logical Aspects of Computational Linguistics (LACL 2014)},
Lecture Notes in Computer Science 8535, Springer, 2014, pp.~57--69.
\newblock \url{https://doi.org/10.1007/978-3-662-43742-1_5}.

\bibitem{ClarkEyraud2007}
A.~Clark and R.~Eyraud,
\newblock Polynomial identification in the limit of substitutable
context-free languages,
\newblock \emph{Journal of Machine Learning Research} 8 (2007), 1725--1745.

\bibitem{Yoshinaka2009}
R.~Yoshinaka,
\newblock Learning mildly context-sensitive languages with multidimensional
substitutability from positive data,
\newblock in \emph{Algorithmic Learning Theory 2009}, LNCS 5809,
278--292.

\bibitem{Yoshinaka2011}
R.~Yoshinaka,
\newblock Efficient learning of multiple context-free languages with
multidimensional substitutability from positive data,
\newblock \emph{Theoretical Computer Science} 412(19) (2011), 1821--1831.

\bibitem{HolzerKonig2004}
M.~Holzer and B.~K{\"o}nig,
\newblock On deterministic finite automata and syntactic monoid size,
\newblock \emph{Theoretical Computer Science} 327(3) (2004), 319--347.
\newblock \url{https://doi.org/10.1016/j.tcs.2004.04.010}.

\bibitem{PaullUnger1959}
M.~C.~Paull and S.~H.~Unger,
\newblock Minimizing the number of states in incompletely specified sequential switching functions,
\newblock \emph{IRE Transactions on Electronic Computers} EC-8(3) (1959), 356--367.
\newblock \url{https://doi.org/10.1109/TEC.1959.5222697}.

\bibitem{Pfleeger1973}
C.~P.~Pfleeger,
\newblock State reduction in incompletely specified finite-state machines,
\newblock \emph{IEEE Transactions on Computers} C-22(12) (1973), 1099--1102.

\bibitem{BarringtonTherien1988}
D.~A.~M.~Barrington and D.~Th{\'e}rien,
\newblock Finite monoids and the fine structure of $\mathsf{NC}^{1}$,
\newblock \emph{Journal of the ACM} 35(4) (1988), 941--952.
\newblock \url{https://doi.org/10.1145/48014.63138}.

\bibitem{GrosshansMcKenzieSegoufin2017}
N.~Grosshans, P.~McKenzie, and L.~Segoufin,
\newblock The power of programs over monoids in $\mathbf{DA}$,
\newblock in: \emph{42nd International Symposium on Mathematical Foundations of Computer Science (MFCS 2017)},
LIPIcs 83, Article 2, pp.~2:1--2:20, 2017.
\newblock \url{https://doi.org/10.4230/LIPIcs.MFCS.2017.2}.


\bibitem{PittWarmuth1993}
L.~Pitt and M.~K.~Warmuth,
\newblock The minimum consistent DFA problem cannot be approximated within any polynomial,
\newblock \emph{Journal of the ACM} 40(1) (1993), 95--142.
\newblock \url{https://doi.org/10.1145/138027.138042}.

\bibitem{ChenFarzanClarkeTsayWang2009}
Y.-F.~Chen, A.~Farzan, E.~M.~Clarke, Y.-K.~Tsay, and B.-Y.~Wang,
\newblock Learning minimal separating DFA's for compositional verification,
\newblock in \emph{Tools and Algorithms for the Construction and Analysis of
Systems (TACAS 2009)}, Lecture Notes in Computer Science 5505,
Springer, 2009, pp.~31--45.
\newblock \url{https://doi.org/10.1007/978-3-642-00768-2_3}.

\bibitem{VazquezPargaGarciaLopez2016}
M.~V\'azquez de Parga, P.~Garc\'ia, and D.~L\'opez,
\newblock A sufficient condition to polynomially compute a minimum separating DFA,
\newblock \emph{Information Sciences} 370--371 (2016), 204--220.
\newblock \url{https://doi.org/10.1016/j.ins.2016.07.053}.

\bibitem{LaumenSnelVaandrager2026}
J.~Laumen, L.~Snel, and F.~Vaandrager,
\newblock An $L^{\#}$ based algorithm for active learning of minimal separating automata,
\newblock in \emph{Computer Aided Verification (CAV 2026)},
Lecture Notes in Computer Science 16683, Springer, 2026, pp.~261--283.
\newblock \url{https://doi.org/10.1007/978-3-032-32526-6_13}.







\bibitem{Howie1995}
J.~M.~Howie,
\newblock \emph{Fundamentals of Semigroup Theory},
\newblock London Mathematical Society Monographs, New Series 12,
Clarendon Press, Oxford, 1995.

\bibitem{RhodesSteinberg2009}
J.~Rhodes and B.~Steinberg,
\newblock \emph{The $q$-Theory of Finite Semigroups},
\newblock Springer Monographs in Mathematics, Springer, New York, 2009.
\newblock \url{https://doi.org/10.1007/b104443}.

\bibitem{Sakarovitch1979}
J.~Sakarovitch,
\newblock Mono\"{\i}des point\'es,
\newblock \emph{Semigroup Forum} 18 (1979), 235--264.
\newblock \url{https://doi.org/10.1007/BF02574188}.

\bibitem{Pin1981}
J.-E.~Pin,
\newblock An analog of the theorem of varieties for cones and cylinders,
\newblock in: \emph{Automata, Languages and Programming}, Lecture Notes in
Computer Science 104, Springer, 1981, pp.~78--90.
\newblock \url{https://doi.org/10.1007/BFb0017300}.

\bibitem{Pin1986}
J.-\'E.~Pin,
\newblock \emph{Varieties of Formal Languages},
\newblock North Oxford Academic, London, 1986.

\bibitem{Almeida1994}
J.~Almeida,
\newblock \emph{Finite Semigroups and Universal Algebra},
\newblock World Scientific, Singapore, 1994.


\bibitem{CostaFlorencioVerwer2014}
C.~Costa Flor\^encio and S.~E.~Verwer,
\newblock Regular inference as vertex coloring,
\newblock \emph{Theoretical Computer Science} 558 (2014), 18--34,
\newblock \url{https://doi.org/10.1016/j.tcs.2014.09.023}.

\bibitem{GareyJohnson1979}
M.~R.~Garey and D.~S.~Johnson,
\newblock \emph{Computers and Intractability: A Guide to the Theory of
NP-Completeness},
\newblock W.~H.~Freeman, San Francisco, 1979.



\bibitem{MargolisPin1984}
S.~W.~Margolis and J.-\'E.~Pin,
\newblock Minimal noncommutative varieties and power varieties,
\newblock \emph{Pacific Journal of Mathematics} 111(1) (1984), 125--135.
\newblock \url{https://doi.org/10.2140/pjm.1984.111.125}.



\bibitem{KuriyamaMCFG2026}
T.~Kuriyama,
\newblock Finite sentence-interface control for learning bounded-fan-out linear MCFGs under fixed monoid typing,
\newblock arXiv:2605.11644, 2026.
\newblock \url{https://arxiv.org/abs/2605.11644}.

\bibitem{Thumm2025}
A.~Thumm,
\newblock Finite semigroups satisfying an identity
\(x_1\cdots x_n\approx\rho(x_1,\ldots,x_n)\),
\newblock arXiv:2509.19216, 2025.
\newblock \url{https://arxiv.org/abs/2509.19216}.

\bibitem{ThummWeiss2026}
A.~Thumm and A.~Wei\ss,
\newblock Efficient compression in semigroups,
\newblock in \emph{43rd International Symposium on Theoretical Aspects of
Computer Science (STACS 2026)}, LIPIcs 364, Article 80, pp.~80:1--80:20,
2026.
\newblock \url{https://doi.org/10.4230/LIPIcs.STACS.2026.80}.

\bibitem{Petrich1969}
M.~Petrich,
\newblock Representations of semigroups and the translational hull of a
regular Rees matrix semigroup,
\newblock \emph{Transactions of the American Mathematical Society}
143 (1969), 303--318.
\newblock \url{https://doi.org/10.1090/S0002-9947-1969-0246984-3}.



\bibitem{Petrich1974}
M.~Petrich,
\newblock The structure of completely regular semigroups,
\newblock \emph{Transactions of the American Mathematical Society}
189 (1974), 211--236.
\newblock \url{https://doi.org/10.1090/S0002-9947-1974-0330331-4}.



\bibitem{BarberRuskucDirectSquare}
C.~Barber and N.~Ru\v{s}kuc,
\newblock Semigroup congruences and subsemigroups of the direct square,
\newblock \emph{Bulletin of the Australian Mathematical Society}
113(1) (2026), 123--134.
\newblock \url{https://doi.org/10.1017/S0004972725100166}.


\bibitem{BarberRuskucDiagonal2026}
C.~Barber and N.~Ru\v{s}kuc,
\newblock Comparing numbers of diagonal subsemigroups and congruences for semigroups,
\newblock arXiv:2602.17286v1, 2026.
\newblock \url{https://arxiv.org/abs/2602.17286}.

\bibitem{ZelinkaTolerance1975}
B.~Zelinka,
\newblock Tolerance in algebraic structures II,
\newblock \emph{Czechoslovak Mathematical Journal} 25 (1975), 175--178.
\newblock MR~0457323; Zbl~0316.08001.

\bibitem{Pondelicek1978}
B.~Pond\v{e}l\'\i\v{c}ek,
\newblock On tolerances on periodic semigroups,
\newblock \emph{Czechoslovak Mathematical Journal} 28(4) (1978), 647--649.
\newblock MR~0498924; Zbl~0399.20053.

\bibitem{Loganathan1986}
M.~Loganathan,
\newblock On tolerance relations,
\newblock \emph{Czechoslovak Mathematical Journal} 36(4) (1986), 662--667.

\bibitem{KumaresanTolerance1984}
K.~Kumaresan,
\newblock A note on the strongly compatible tolerances on an arbitrary semigroup,
\newblock \emph{Czechoslovak Mathematical Journal} 34(2) (1984), 319--321.
\newblock \url{https://eudml.org/doc/13454}.

\end{thebibliography}

\end{document}
